%  Guide de l'usager du logiciel LatMRG
%
\documentstyle[twoside,11pt,myarticle]{article}
% \usepackage{url}
% \input url.sty
\def\url#1{{\tt #1}} 

% \includeonly {TESTCONF}
\mytwoheads
%%\dateheadtrue

\def\ray#1    {\fbox {\footnote {\ }}\ \footnotetext { From Raymond: #1}}
\def\pierre#1 {\fbox {\footnote {\ }}\ \footnotetext { From Pierre: #1}}
\def\hray#1 {}
\def\hpierre#1 {}
% Changing  \ray  to  \hray  will temporarily hide the comment.
% Comments can then be "hiddenly" saved for further reference.

\def\title#1    {\vfill\begin {center} \Large\rm #1 \end {center} \vfill}
\def\author#1   {\begin {center} \large\rm #1 \end {center} \medskip}
\def\abstract   {\section* {ABSTRACT} \null }
\def\keywords#1 {\noindent #1 \vfill\vfill}
\def\thanks     {\section* {Acknowledgements}}
%\def\refer      {\begin {thebibliography}{99}
%\def\endrefer   {\end {thebibliography}}
\def\newblock   {\hskip .11em plus .33em minus .07em} 

\def\?{\discretionary{}{}{}}  % Same as \- but does not print the - sign
                              % (useful to hyphenate in math mode).
\def\To {\Rightarrow}
\def\eqdef {\buildrel \rm def \over =}    % Egal par definition.
\def\tod {\buildrel \rm d \over \To}
\def\iid {{\rm i.i.d.}}
\def\mod{{\rm\ mod\ }}        % Modulo.
\def\MOD{{\rm\ MOD\ }}
\def\div{{\rm\ div\ }}        % Division entiere.
\def\DIV{{\rm\ DIV\ }}
\def\var{{\rm\ var }}        % Variance.
\def\Var{{\rm\ Var }}
\def\arg{{\rm\ arg }}          % Argument.
\def\dim{{\rm dim }}
\def\cA {{\cal A}}
\def\cB {{\cal B}}
\def\cH {{\cal H}}
\def\cI {{\cal I}}
\def\cJ {{\cal J}}
\def\cF {{\cal F}}
\def\cL {{\cal L}}
\def\cN {{\cal N}}
\def\cP {{\cal P}}
\def\Cov{{\rm\ Cov }}
\def\Is {[0,1)^s}
\def\fFour{\hat{f}({\bf h})}
\def\hdual{{\bf 0} \neq{\bf h} \in L_s^{*} }

\def\o   {{(o)}}
\def\a   {{*}}
\def\ba  {{\bf a}}
\def\bA  {{\bf A}}
\def\bb  {{\bf b}}
\def\bB  {{\bf B}}
\def\bc  {{\bf c}}
\def\be  {{\bf e}}
\def\bI  {{\bf I}}
\def\bh  {{\bf h}}
\def\bk  {{\bf k}}
\def\bm  {{\bf m}}
\def\bM  {{\bf M}}
\def\bT  {{\bf T}}
\def\bu  {{\bf u}}
\def\bU  {{\bf U}}
\def\bv  {{\bf v}}
\def\bV  {{\bf V}}
\def\bw  {{\bf w}}
\def\bW  {{\bf W}}
\def\bx  {{\bf x}}
\def\bX  {{\bf X}}
\def\by  {{\bf y}}
\def\bz  {{\bf z}}

\def\rmB  {{\rm B}}
\def\rmL  {{\rm L}}
\def\rmM  {{\rm M}}
\def\rmP  {{\rm P}}
\def\rmR  {{\rm R}}
\def\rmZ  {{\rm Z}}

\def\tbA {\tilde{\bf A}}
\def\tbs {\tilde{\bf s}}
\def\tbu {\tilde{\bf u}}
\def\bzero {{\bf 0}}
\def\es{\hspace{0.5mm}e\hspace{0.6mm}}
\def\card {\mbox{\rm card}}

\def\Re{{\rm I\kern-0.2em R}} % Ensemble des nombres reels.
\def\R{\Re}
\def\N{{\rm I\kern-0.2em N}}  % Ensemble des nombres naturels.
\def\F{{\rm I\kern-0.2em F}}  % Ensemble des nombres naturels.
\def\Z{{\sf Z\kern-0.4em Z}}  % Ensemble des nombres entiers.
\def\B{{\cal B}}
\def\square{\ $\Box$}
\def\eqdef {\buildrel \rm def \over =}    % Egal par definition.
\def\<{$\langle$}        %  Signe <
\def\>{$\rangle$}        %  Signe >
\def\q{$\kern1.4em$}     % Space for indentation in slides and programs.


%%%%%%%%%%%%%%%%%%%%%%%%%%%%%%%%%%%%%%%%%%%%%%%%%%%%%%%%%%%%%%%%%%%%%
\begin {document}

\begin  {titlepage}
\null
\begin{flushright}
 {\it  Version: November 2000 } 
\end{flushright}
\vfill
%
\title  { \huge\bf  LatMRG \medskip \\ \LARGE\bf User's Guide\\ \bigskip
          \large\bf A Modula-2 software for the theoretical analysis \\
            of linear congruential and multiple recursive \\
            random number generators }

\author { Pierre L'Ecuyer \ and \ Raymond Couture \\ \medskip
          \small\rm  % GERAD \\ and \\
          D\'epartement d'Informatique et de Recherche Op\'erationnelle \\
          Universit\'e de Montr\'eal, C.P. 6128, Succ.\ Centre-Ville,
          Montr\'eal, Canada, H3C 3J7 }

{\noindent\large\bf Abstract. \ }
{\tt LatMRG} is a software toolkit for examining theoretical properties of
linear congruential or multiple recursive random number generators.
It is implemented as a library of modules written in Modula-2.
It offers tools to check whether a generator has maximal period or not,
to apply the lattice and spectral tests (in large dimensions),
and to perform computer searchs for good (or bad) generators 
according to different criteria.
One can analyse the lattice structure of points formed by successive
values in the generator's sequence, or formed by ``leapfrog'' values.
Generators with large moduli and multipliers (e.g. numbers of many
hundreds of bits), as well as combined generators, can also be analyzed.
Multiply-with-carry generators can also be studied by analyzing their
corresponding linear congruential generators. 

\keywords {{\large\bf Keywords:} 
            Random number generation; linear congruential generators;
            multiple recursive generators;  multiply-with-carry;
            lattice structure; spectral test; maximal period }

\end {titlepage}

\pagenumbering{roman}
\tableofcontents
%\include{avantpro}
\newpage
%
%                   Background, overview, examples, guidelines.
\pagenumbering{arabic}
\input sec1.tex
\input sec2.tex
\input sec3.tex
%
%                   Programs in executable form.
\appendix
\typeout {Appendix A  ***}
\thispagestyle {nomark}
\addcontentsline{toc}{part}{APPENDICES}
\addcontentsline{toc}{section}{\protect\numberline{A.}{
  Programs in Executable Form }}
\null\vfill\vfill
\centerline {\large\bf APPENDIX A}
\vfill
\centerline {\large\bf Programs in Executable Form}
\vfill\vfill
%\eject\null\vfill         % Page blanche.

\defmodule {maxper}

This program verifies whether a given generator has maximal period or not.
Integers are represented using the {\tt SuperInteger} type.
The modulus $m$ must be prime.
To verify the maximal period conditions, the factorizations of $m-1$ and
$r = (m^k-1)/(m-1)$ are required.  They can be found by the program
or provided by the user in a file.
We warn that factoring $r$ can take a {\em huge} amount of time.

The data for {\tt maxper} must be placed in a file with extension
``{\tt .dat}'', according to the format displayed in Figure~\ref{maxper.dat}.
The data fields have the same meaning as for {\tt latl}.
To run the program, type ``{\tt maxper \em\<file\>}'', 
where {\em\<file\>} is the name of the data file, without extension.
Currently, only {\tt LCG} and {\tt MRG} are allowed for {\em TypeGen}.

%%  DATA FILE  %%%%%%%%%%%%%%%%%%%%%%%%%%%%%%%%%%%%%%%%
\setbox0=\vbox {\hsize=12.0cm \setverbatim\it\baselineskip=13pt

           TypeGen
           $m$
           $k$
           $F1$ $[$\<file1\>$]$
           $F2$ $[$\<file2\>$]$
           $a_1$
           \vdots
           $a_k$

}
\begin{figure}[h] \centerline {\boxit{\box0}}
\caption { Data file format for {\tt maxper}. }  \label {maxper.dat}
\end{figure}


\defmodule {findmk}

This program is an interface to the module {\tt PRIM}.
It searches and prints the list of all prime integers $m$ such
that $2^e+ c_1 \le m\le 2^e + c_2$ and $r = (m^k-1)/(m-1)$ is prime.
We may also request that $(m-1)/2$ be also prime, by setting
the boolean variable {\em Safe\/} to {\tt TRUE}.

The program reads the data and calls the procedure {\tt FindPrimesmr}
in the module {\tt PRIM}, with the corresponding parameters.
The data for {\tt findmk} must be placed in a file with extension
``{\tt .dat}'', in the format of Figure~\ref{findmk.dat}.
To run the program, type ``{\tt findmk \em\<file\>}'', 
where {\em\<file\>} is the name of the data file, without extension.
The results will be written in {\em\<file\>}{\tt .res}.

%%  DATA FILE  %%%%%%%%%%%%%%%%%%%%%%%%%%%%%%%%%%%%%%%%
\setbox1=\vbox {\hsize=12.0cm \setverbatim\it\baselineskip=13pt

           $k$
           $e$
           $c_1$
           $c_2$
           Safe

}
\begin{figure}[h] \centerline {\boxit{\box1}}
\caption { Data file format for {\tt findmk}. }   \label {findmk.dat}
\end{figure}


\defmodule {latl}

This program applies the spectral and/or lattice test(s) to a given
generator.  It uses the {\tt LONGINT} representation for the multipliers
and modulus (module {\tt MULTLI}) and the {\tt LONGREAL} representation
for the bases (module {\tt BASISLR}).
The generator analyzed can be a combined generator with $J$ components,
and expressed that way in the data file.
\iffalse  %%%%%%
If $J=1$ and $m$ is a power of two expressed as
$(x\ y\ z) = (2\ y\ 0)$ in the data file,
then the lattice associated to the multiplicative LCG with power
of two modulus, as described in \cite{FIS90}, is analyzed.
Otherwise, the lattice analyzed is the smallest one that contains all
the points, for all possible starting values (i.e. the union of all subcycles,
when there are subcycles).
\fi  %%%%%%
To run the program, type ``{\tt latl \<\em file\>}'', where \<{\em file}\> is
the name of the data file, without extension.

The data must be in a file with extension
``{\tt .dat}'', according to the format of Figure~\ref{lat.dat}.
The data fields have the same meaning as for {\tt seekl},
except that several fields have been removed, 
$a_i$ replaces $b_i$ and $c_i$,
there is a single category for the dimensions, 
and the options for {\em LatticeInfo} and {\em ResultForm\/} are different.
\begin {description}
\item [$(a_1,\ldots,a_k)$ :] is the vector of (integer) multipliers.
  They can be given in format (a) or (b) described in {\tt seekl}.
\item [{\em LatticeInfo} ~\<{\em Norm\/}\> :]
  If the value of {\em LatticeInfo} is {\tt Spectral} 
  or {\tt BeyerSpectral} or {\tt SpectralP},
  then the field {\em Norm\/} must appear.
  It indicates which type of normalization is used
  in the definition of $S_t$.
  The allowed values are {\tt (BestLat, Laminated, Rogers, Minkowski)}.
\item [{\em ResultForm} :]
  Selects in which form the results will be given.
  The possible values are {\tt (Terminal, RES, TEX)}.
  Lowercases are also allowed.
  \begin {description}
   \item [{\tt Terminal}] indicates that the results will be given on the
     terminal screen.
   \item [{\tt RES}] says that the results will be in a file with the same
     name as the data file, but with extension ``{\tt .res}''.
   \item [{\tt TEX}] asks the program to produce a file intended for
     \LaTeX, with extension ``{\tt .tex}''.
  \end {description}
\end {description}

%%  DATA FILE  %%%%%%%%%%%%%%%%%%%%%%%%%%%%%%%%%%%%%%%%
\setbox0=\vbox {\hsize=12.0cm \it \begin {tabular}{l}
           ReadGenFile ~$[$\<file0\>$]$ \\
           $J$ \\
           $\left.\begin {array}{l}
           \mbox {TypeGen} \\
           m \\
           k \\
           a_1 \\
           \vdots \\
           a_k \\
           \end {array} \ \ \ \right\} \mbox { (Repeat $J$ times) }$ \\
           Dim$(0)$ ~Dim$(1)$ \\
           LatticeInfo  ~\<{\em Norm\/}\> \\
           LatticeType \\
           LaGroupSizes ~Spacing\\
           VerifyBB \\
           MaxNodesBB \\
           ResultForm \\
\end {tabular}
}
\begin{figure}[htb] \centerline {\boxit{\box0}}
\caption { Data file format for latl. }  \label {lat.dat}
\end{figure}

\defmodule {seekl}

(To be rewritten....)

This program performs a search for the ``best'' (or ``worst'') multiple
recursive generators or multiply-with-carry generators of a given
form, based on one of the criteria $Q_T$ or $M_T$ defined in Section 1.3
It produces a report listing the retained generators, their properties,
and various statistics on the search.

The set of dimensions in which the test is applied can be partitioned
into a certain number of intervals, or categories,
and one can use a different selection criterion for the generators within
each category.  One can also impose bounds on the figure of merit
within each category.  See the data fields for $C$, {\it MinMerit}
and {\it MaxMerit\/} below.
For example, one can consider only the generators with $M_8 \ge 0.6$,
and among these, retain the list of generators with the {\em smallest\/}
value of $M_{12}$.
As another example, one can retain the 2 generators with the highest $M_8$,
the 2 generators with the highest $M_{16}$ and the eight generators
with the highest $M_{32}$.

One can search for combined MRGs with $J$ components, or simple MRGs $(J=1)$,
or multiply-with-carry (MWC) generators.
For a MWC, one simply analyzes the corresponding LCG, which is a special
case of an MRG.  Therefore, in what follows, we use the term
`MRG component' to denote either and MRG or a MWC.
For a simple MRG (or for each component, in case $J > 1$), with
given modulus $m$ and order $k$, the program searches for vectors
of multipliers inside the
region bounded by the vectors $b = (b_1,\dots,b_k)$ and $c = (c_1,\dots,c_k)$
such that $-m < b_i\le c_i < m$ for each $i$.
The search can be exhaustive in that region, or random.
One can search only among maximal period generators (for each component),
or not consider the period and examine only the lattice structure.
The former (checking maximal period conditions) can be done only
if $m$ is prime, or if $k=1$ and $m$ is a power of a prime.
The program can also list the retained generators in a file, in a format
more compact than for the result file, and can re-use that file as input
to the program, in a later run.
This could be useful, for example, if one wishes to perform first a screening
over a large region, based on a criterion that does not require expensive
computations, and then do a second pass over the retained generators,
based on a more stringent criterion, such as looking at the lattice structure in
higher dimensions, and/or verifying the results by performing all computations
using error bounds.


%%%%%%%%%%%%%%%%%%%%%%%%%%%%%%%%%%%%%%%%%%%%%%%%%%%%%%%%%%%%
\subsection* {Method of search}

For an exhaustive search for MRGs, all vectors of multipliers
of the form $a = (a_1,\dots,a_k)$
such that $b_i\le a_i\le c_i$ for $i=1,\dots,k$ will be examined, for
a total of $N_v = \prod_{i=1}^k (c_i-b_i+1)$ vectors.
This holds for each component.
Therefore, if there are $J$ components and $N_{v,j}$ vectors are
examined for component $j$, then a total of $\prod_{j=1}^J N_{v,j}$
generators are examined.

For a random search for MRGs, we fix a number of subregions (clusters)
we want to examine, and the size $h_i$ of each subregion in dimension $i$,
for $i=1,\dots,k$.  The program will examine a total of
$n\prod_{i=1}^k h_i$ vectors of multipliers (for each MRG component)
by repeating $n$ times the following:
% \begin {verse}
  For $i=1,\dots,k$, generate $\alpha_i$ randomly, uniformly over the set
  $\{b_i,\dots,c_i-h_i+1\}$;  then, examine all the vectors
  $a = (a_1,\dots,a_k)$ such that $\alpha_i\le a_i\le \alpha_i+h_i-1$
  for each $i$.
% \end {verse}

When examining a vector $a$, the program first checks if the maximal
period conditions are satisfied, if this is required.
For prime modulus $m$, the condition (a) (see Section 1.2) is verified
only once for each distinct value of $a_k$ (which corresponds to
$\prod_{i=1}^{k-1} h_i$ different vectors).
To verify the maximal period conditions, the factorizations of $m-1$ and
$r = (m^k-1)/(m-1)$ are required.  They can be found by the program,
if desired, or provided by the user in a file (see below).
We recall that factoring $r$ can take {\em huge} amounts of time.
So, avoid redoing the factorization unnecessarily.
For the MRGs, these factorizations are necessary only when $m$ is
prime and maximal period is required.

If $a$ is not rejected by the maximal period test, then we move forward
to the next MRG component and try all the vectors for that next component
(by exhaustive or random search) and examine their combination with
the currently examined multipliers for the previous components.
For each combined generator, the values of $d_t$ and/or $q_t$ are 
computed for dimensions $k+1,\dots,T$.  

The program always keeps lower and upper bounds
on the figure of merit ($M_T$ or $Q_T$), in each dimension,
for the generator to be worth considering.  
The initial values of these bounds are given by the user
in the fields {\em MinMerit} and {\em MaxMerit}.
The lower bound can be 0.0 and the upper bound can be 1.0, 
which means that there can be no effective bounds for some
categories if desired.

As soon as a generator has a figure of merit below the lower
bound in a given dimension, or above the upper bound for its
category (after the computations for all the dimensions in this
category have been completed), then this generator is immediately
discarded and no further computations are made for it.
This can save enormous amounts of time in the case of very large
searches up to high dimensions, because with good bounds,
few generators will reach the large dimensions.

During execution, only the bounds for the last category can be
modified.  If the figure of merit for the last category is to be
{\em maximized\/}, when we have found enough (i.e., {\em NbGen($C$)})
generators with a figure of merit $\ge \sigma$, where $\sigma$
is larger than the lower bound for the last category, then we raise
this lower bound to $\sigma$.
Similarly, if we minimize in the last category, we can lower 
the upper bound when we have enough generators beating the bound.
In the case where the figure of merit is to be {\em maximized in 
all\/} the categories, then a generator is also discarded
as soon as its figure of merit in {\em any\/} dimension gets below 
the lower bound of the last category.

In the case where the selection criterion is {\tt Spectral} or
{\tt SpectralH}, and if $T>8$, the distances between
hyperplanes are computed for dimensions up to $T$, but the selection of
generators is based only on $M_{8}$, because the Hermite constants
defining $d_t^*$ in that case are known only for $t\le 8$.

The execution (cpu) time is checked before testing each new generator.
When it exceeds the cpu time limit given in the data file,
the search is aborted and the partial results are printed.

%%  DATA FILE  %%%%%%%%%%%%%%%%%%%%%%%%%%%%%%%%%%%%%%%%
\setbox0=\vbox {\hsize=12.0cm \it \begin {tabular}{l}
           ReadGenFile ~$[$\<file0\>$]$ \\
           $J$ \\
           $\left.\begin {array}{l}
           \mbox {TypeGen} \\
           m \\
           k \\
           w \\
           \mbox {PerMax} \\
           \mbox {ImplemCond} ~[\mbox{\<NbMaxBits ~HighestBit\>}] \\
           F1 ~[\mbox{\<file1\>}] \\
           F2 ~[\mbox{\<file2\>}] \\
           b_1 \\
           c_1 \\
           \vdots \\
           b_k \\
           c_k \\
           \mbox {SearchMethod $[$n H $H_k]$ } \\
           \end {array} \ \ \ \right\} \mbox { (Repeat $J$ times) }$ \\
           $C$ \\
           NbGen$(1)$ MinMerit$(1)$ MaxMerit$(1)$ 
                $d_1$ $t_{11}$ $\cdots$ $t_{1,d_1}$ \\
           \vdots \\
           NbGen$(C)$ MinMerit$(C)$ MaxMerit$(C)$ 
                $d_C$ $t_{C1}$ $\cdots$ $t_{C,d_C}$ \\
           Criterion  ~\<{\em Invert\/}\> ~\<{\em Norm\/}\> \\
           LatticeInfo \\
           LatticeType \\
           LaGroupSizes ~Spacing\\
           VerifyBB \\
           MaxNodesBB \\
           TimeLimit \\
           $S1 ~S2$ \\
           ResultForm \\
\end {tabular}
}
\begin{figure}[htb] \centerline {\boxit{\box0}}
\caption { Data file format for {\tt seekl}. }  \label {seekl.dat}
\end{figure}

%%%%%%%%%%%%%%%%%%%%%%%%%%%%%%%%%%%%%%%%%%%%%%%%%%%%
\subsection* {The data file}

The data for {\tt seekl} must be placed in a file with extension
``{\tt .dat}'', according to the format displayed on Figure~\ref{seekl.dat}.
The fields in square brackets are optional (depending on the value taken
by the first field on the line).
The meaning of all data fields is explained below.
To run the program, type ``{\tt seekl \<\em file}\>'', where \<{\em file}\> is
the name of the data file, without extension.
The results will be in a file with the same name, with extension
``{\tt .res}'' or ``{\tt .gen}'' (see {\em ResultForm\/}).

Comments may be inserted after data, on the same line,
separated from the data by at least one blank. 
Moreover, any line starting with ``\%'' or ``\#'' is considered
as a comment.

The values of $m$, $b_i$, and $c_i$ in the data file can be given in one
of the two following formats:
\vspace {-10pt}
\begin {itemize}
\item [a)]  An integer giving the value directly, in base 10.
  In this case, there {\em must\/} be some other non-numeric text 
  (e.g., a comment) on this data line after the integer.
\item [b)]  Three integers $x$, $y$, and $z$ on the same line, separated
  by at least one blank.  The retained value will be $x^y+z$ if $x\ge0$, and
  $-((-x)^y+z)$ if $x<0$.  The value of $y$ must be positive.
  For example, $(x\ y\ z) = (2\ 5\ -1)$ will give 31, while
  $(x\ y\ z) = (-2\ 5\ -1)$ will give $-31$ (not $-33$).
\end {itemize}
For the program {\tt seekl}, all these numbers must fit in a {\tt LONGINT}.
For larger numbers, one must use the program {\tt seeks}.

%%%%%%%%%%%%%%%%%%%%%%%%%%%%%%%%%%%%%%%%%%%%%%%%%%%%%%%%%%%%%%%
\subsection* {Meaning of the data fields}

\begin {description}
%
\item [{\em ReadGenFile} and \<{\em file0}\> :]
  {\tt BOOLEAN} and file name (without extension).
  When {\em ReadGenFile} is {\tt FALSE},
  the search is made according to the values of the fields below.
  When {\tt TRUE}, the generators to be looked at are those
  listed in the file \<{\em file0}\>{\tt .gen}.
  This must be a file of type ``{\tt .gen}'', produced by this program
  with the {\tt GEN} option for the {\em ResultForm} data field.
  In that case, only those generators listed in that file are examined
  and the vectors $b$ and $c$ below are not used.
%
\item [{\em J} :]
  Number of components in the combined generators.
  Must be an integer in the range $[1..{\tt MaxJ}]$.
  When $J > 1$, we look for combined generators.
%
\item [{\em TypeGen} :] Can be:
  \begin {description}
  \item [{\tt MRG}] means that this component is an MRG.
  \item [{\tt MWC}] means that this component is a multiply-with-carry 
    (MWC) generator.
    Each MWC generator is converted by the program to its
    corresponding LCG (see, e.g., \cite{rCOU97a,rLEC98a}).
  \item [{\tt MMRG}] means a matrix MRG.
  \end {description}
%
\item [{\em m, k}, \<{\em w}\> :]
  Modulus, order of the recurrence, and size of the matrix-type 
  coefficients.
  Must be positive integers, with $k<$ {\tt MMaxDim}.
  For $m^k >= 2^{31}$, use the program {\tt seeks} instead of {\tt seekl}.
  The value of $w$ should be specified only for MMRG generators.
  Otherwise it is assumed to be 1.
%
\item [{\em Permax} :]
  {\tt BOOLEAN} variable.  
  {\tt TRUE} if maximal period is required, {\tt FALSE} otherwise.
  When set to {\tt TRUE}, $m$ must be expressed in the data file in the
  form (b): $(x\ y\ z)$, otherwise {\em Permax\/} will be put back to
  {\tt FALSE}.  The software assumes that $m$ is prime, 
  unless $z=0$ and $y > 1$, in which case it assumes that $x$ is prime.
  In the latter case, one must have $k=1$, otherwise {\em Permax\/}
  will be set back to {\tt FALSE}.
%
\item [{\em ImplemCond} and \<{\em NbPow2 ~HighestBit}\> :]
  Can be {\tt NoCond}, {\tt AppFact}, or {\tt SumPow2}.
  If {\tt NoCond}, then no conditions are imposed on the
  multipliers $a_i$.
  If {\tt AppFact}, then the multipliers must satisfy the 
  ``approximate factoring'' condition $a_i (m\mod a_i) < m$ for each $i$.
  MRGs are usually easier to implement under this condition \cite{rLEC90a}.
  If {\tt SumPow2}, then the positive integers {\em NbPow2}
  and {\em HighestBit\/} must appear and they indicate that 
  for each $i$, the multiplier $a_i$ must be the sum of at most 
  {\em NbPow2} (positive or negative) powers of 2, 
  with the highest power of 2 not exceeding $2^e$ in absolute value,
  where $e =$ {\em HighestBit}.  
  For example, if {\em SumPow2\/} = 2 and {\em HighestBit\/} = 30,
  there are $30 \times 31 /2$ possibilities for choosing the 2 powers of 
  2 and 4 possibilities for choosing their signs, yielding
  1860 cases where $a_i$ is obtained from exactly 2 powers of 2.
  If one adds the $62$ cases where $a_i$ is $\pm$ a power of 2,
  this gives a total of 1922 possibilities for $a_i$.
%  In the case of {\tt SumPow2}, the bounds
%  $b_j$ and $c_j$ on the multipliers (see below) are automatically
%  (implicitly) reset to $b_j = -(2^{e+1}-1)$ and $c_j = 2^{e+1}-1$.
%
\item [$F1$ and \<{\em file1}\> :]
  This line of data (and also the following one) is used only if 
  maximal period is required and $m$ is assumed to be prime
  (see the {\em Permax\/} field).  Otherwise, the program just skips it
  (but the line must be there).
  $F1$ indicates how the factors of $m-1$ are to be found and
  \<{\em file1}\> is a file name.
  The values allowed for $F1$ are {\tt (Decomp, DecompWrite, Read)}.
  \begin {description}
  \item  [{\tt Decomp}] means that the program itself will factorize
    $m-1$.  In this case, the field \<{\em file1}\> is not used and can be
    omitted.
    To factorize, the program uses the procedure
    {\tt Factorize} in the module {\tt SUPFACT} of the package SENTIERS
    \cite {iLEC88d}, with no CPU time limit.
    It is the responsibility of the user to make sure that the
    factorization will take a reasonable time.
  \item [{\tt DecompWrite}] means the same as {\tt Decomp}, except that
    the program will also write the factors (larger than 1) in the file
    \<{\em file1}\>, one factor per line.
  \item [{\tt Read}] indicates that $m-1$ is already factorized and that
    the factors will be read from file \<{\em file1}\>, in the same format.
    The factors need not be sorted, but must be one per line, and repeated
    factors must be on successive lines.  The factorization must be
    complete and the program will check if the product of all the factors
    is really equal to $m-1$.
  \end {description}
%
\item [$F2$ and \<{\em file2}\> :]
  This data line is similar to the previous one, except that it concerns
  $r=(m^{kw}-1)/(m-1)$ instead of $m-1$.
  In this case, it is possible that $r$ be prime when $m$ is prime
  (in contrast to $m-1$, which is then even).  
  Therefore, the additional value {\tt Prime} is
  allowed for $F2$, so that the set of possible values is
  {\tt (Decomp, DecompWrite, Read, Prime)}.
  The module {\tt SUPFACT} in SENTIERS \cite{iLEC88d} can be used independently
  to factorize $r$ or to check its primality.
  \begin {description}
  \item [{\tt Prime}] indicates that $r$ is prime.
  \end {description}
%
\item [$b=(b_1,\ldots,b_k)$ \rm and $c=(c_1,\ldots,c_k)$~:]
  The $b_i$ and $c_i$ are integers such that
  $-m < b_i \le c_i < m$ for $i=1,\dots,k$.
  They determine the boundary of the (rectangular) area of search.
%  When {\em ImplemCond} = {\tt SumPow2}, these values are not used.
  For MMRGs ($w>1$), each $b_i$ and $c_i$ is replaced by a matrix of
  size $w\times w$.  Each of these matrices is given line by line,
  $w$ entries on each line of the data file.
%
\item [{\em SearchMethod} and $n,\, H,\, H_k$ :]
  {\em SearchMethod} can be {\tt Exhaust} or {\tt Random}.
  \begin {description}
  \item  [{\tt Exhaust}] means that the search will be exhaustive over
    all the region determined by $b$ and $c$.
    The other parameters of this line are then useless and can be omitted.
  \item  [{\tt Random}] asks for a random search, performed as described
    in the first subsection of the description of {\tt seekl}.
    The integer $n$ gives the number of subregions (clusters) to examine. 
    If $w=1$, $H$ determines the size of these subregions, except for
    the $k$th element of the vector $h$, where $H_k$ determines the size.
%    The parameter $Hk$ is not used if $k = 1$.
    The vector $h=(h_1,\ldots,h_k)$ is computed by the program as follows:
    $h_i = \min (H, c_i-b_i+1),\ i=1,\ldots,k-1$, and
    $h_k = \min (H_k, c_k-b_k+1)$.
    If $w > 1$, ?????
  \end {description}
%
\item [$C$:] 
  The program will retain $C$ lists (or categories) of generators, 
  according to the specifications given below.
  One must have $1 \le C \le {\tt MaxCat}$.
%
\item [{\em NbGen$(1)$ ~MinMerit$(1)$ ~MaxMerit$(1)$} 
        $d_1$ $t_{11}$ $\cdots$ $t_{1,d_1}$]
\item [ \hskip 1in\vdots ] \
\item [{\em NbGen$(C)$ ~MinMerit$(C)$ ~MaxMerit$(C)$} 
        $d_C$ $t_{C1}$ $\cdots$ $t_{C,d_C}$] 
  \ \\ \ \\
  Integers which must satisfy the constraints: 
  \begin {eqnarray*}
   & 0\le {\em NbGen}(i) \le {\tt MaxNbGen},\\
   & 0.0 \le {\it MinMerit}(c) \le {\it MaxMerit}(c) \le 1.0 &
            \qquad\mbox{ for all } c, \\
   & {\it MinMerit}(c) \le {\it MaxMerit}(c-1) & \qquad\mbox{ for } c > 1,\\
   & d_1 \le \cdots\le d_C, \\
   & {\tt BMaxDim} \ge t_{c1} \ge \cdots\ge t_{c d_c}\ge d_c &
            \qquad\mbox{ for all } c, \\ 
% t_{1d} \le \cdots\le t_{C d} \mbox{ for all } d. 
%  -> Try keeping this constraint out.
  \end {eqnarray*}
  where {\tt BMaxDim} is defined in the module {\tt CONFIG}.

  There will be {\em NbGen$(c)$} generators in the $c$th list.
  The real numbers {\em MinMerit$(c)$ and MaxMerit$(c)$} represent 
  the minimal and maximal values of the figure of
  merit to keep a generator in the $c$th list.
  That is, only the generators that satisfy 
 \eq  
  \mbox{\it MinMerit}(c) \le\sigma\le \mbox{\it MaxMerit}(c) 
                                                   \label{con-merit}
 \endeq 
  are considered for category $c$ and for the categories $c' > c$.  
  Because of this, {\it MinMerit}$(c)$ must be nondecreasing in $c$,
  because the lower bound for category $c$ also
  apply categories $c' > c$.
  As soon as a generator does not satisfy the constraint (\ref{con-merit})
  in a given dimension for a given category $c$,
  it is discarded and no more time is spent to test it for $c' > c$.
  When looking for good generators, one normally sets 
  {\em MaxMerit}$(c)$ to 1.0 for each $c$.

  For category $c$, the test (Beyer, spectral, \dots, depending on
  the criterion) is performed in dimensions $kw+1$ to $t_{c1}$ 
  for the vectors of successive values,
  the pairs $(u_1, u_{j})$ are tested for $1 <j\le t_{c2}$,
  the triples $(u_1, u_{j}, u_{\ell})$ are tested for $1 <j <\ell\le t_{c3}$,
  and so on.  
  If lacunary indices are specified in the field {\em LaGroupSizes},
  all of this is applied to the vectors obtained {\em after\/} taking
  the lacunary indices, and the test for successive values starts in 
  2 dimensions instead of $kw+1$.
  To {\em minimize\/} (instead of maximizing) the criterion within a 
  given category $c$, place a {\em negative sign\/} in front of the value
  of $d_c$ in the data file.
  The program will then retain the {\em worst\/} instead of the
  {\em best\/} generators with respect to the category $c$.
%
\item [{\em Criterion} ~\<{\em Invert\/}\> ~\<{\em Norm\/}\> :]
  Specifies the merit criterion for ranking the generators for each category.
  The field {\em Invert\/} can be either {\tt Invert} or left blank.
  In the former case, the value of the criterion is
  replaced by its multiplicative inverse in the results.
  The admissible values of {\em Criterion\/} are:
  \begin {description}
  \item [{\tt Beyer}] means that the criterion is $Q_T$.
    The program will retain the generators with the largest 
    (or smallest) $Q_T$ in each category.
  \item [{\tt Spectral}] means that the criterion is $M_T$.
    The program will retain the generators with the largest values
    (or smallest) $M_T$ in each category.
  \item [{\tt SpectralP}] is similar to {\tt Spectral}, except that
    it is based on the shortest vector in the {\em primal\/} lattice.
  \item [{\tt SpectralL1}] is similar to {\tt Spectral}, except that
    the length of the (dual) vectors are measured with the $L_1$ norm,
    minus 1.  The length of the shortest dual vector is then an upper
    bound on the minimal number of hyperplanes that cover all the points
    (\cite{rDIE75a} and \cite{rKNU97a}, Exercises 3.3.4-15 and 16).
  \item [{\tt Palpha} {\em Alpha}] means that the criterion is $P_\alpha$,
    with the value of $\alpha$ specified just after the keyword {\tt Palpha}.
    Specialized (much faster) implementations have been made for $\alpha = 2, 4, 6, 8$.
  \end {description}
  If {\em Criterion\/} is {\tt Spectral} or {\tt SpectralP}, 
  then {\em Norm\/}
  must appear and it indicates which type of normalization is used
  in the definition of $S_t$.  The admissible values are
  \begin {description}
  \item [{\tt Hermite}] means that we use Hermite's bound.
    In this case, $S_t$ is computed only up to dimension 8
    and the criterion is $M_{\min(T, 8)}$.  
    This is because for $t>8$, $d_t^*$ is
    unknown, so $M_t$ cannot be computed.
  \item  [{\tt Laminated}] means that we use for $d_t^*$ the value
    of $d_t$ that corresponds to the laminated lattice in dimension $t$.
    The criterion is then $M_{\min(T, 48)}$ 
    (the programs knows this $d_t^*$ only for $t\le 48$).
  \item  [{\tt Rogers}] means that $d_t^*$ is obtained from Rogers'
    bound.  The criterion is then $M_{\min(T, 48)}$.
  \end {description}
  If {\em Criterion\/} is {\tt Palpha}, the field {\em Norm} must also
  appear.  The admissible values for {\tt Palpha} are:
  \begin {description}
%  \item [{\tt ????? }] (Check the paper with Hickernell) ?????
  \item [{\tt Weights}] followed on the next line by positive real numbers 
  $\beta_0,\dots,\beta_t$ where $t= t_{C,1}$.  This means that the criterion
  will be $\tilde P_{\alpha}$ defined in (\ref{eq:Palpha}).
  In this case, $\alpha$ must be even.  The $\beta_j$'s are the weights
  given to the different dimensions in the weighted version of $P_{\alpha}$.
  \end {description}
%
\item [{\em LatticeInfo} :]
  Indicates which figures of merit we want to compute and print.
  The possibilities are the same as for the criterion, but we can compute
  and print more than one.  It suffices to write more than one keywords
  on this line.
  Note that $q_t$ typically requires much more time to compute than $d_t$.
  \begin {description}
  \item [{\tt Beyer} :] Prints the values of $q_t$ up to dimension $T$.
  \item [{\tt Spectral} :] Prints the values of $d_t$ up to dimension $T$ and
    $S_t$ for $t\le\min (T, T^*)$, where $T^*=8$ if {\em Norm} = 
    {\tt Hermite} and $T^* = 48$ otherwise.
  \item [{\tt SpectralP}] is similar to {\tt Spectral}, except that
    it is based on the shortest vector in the {\em primal\/} lattice.
  \item [{\tt SpectralL1}] is similar to {\tt Spectral}, except that
    the length of the (dual) vectors are measured with the $L_1$ norm,
    minus 1.  The normalization is based on Minkowski's bound
    (\cite{sFIS96a}, p. 618).
  \item [{\tt Palpha} {\em Alpha} :] prints $P_\alpha$
    with the value of $\alpha$ specified.
  \end {description}
%
\item [{\em LatticeType} :]
  Indicates whether to analyze the lattice generated by all possible states,
  or a sublattice generated by the set of recurrent states or by a
  subcycle of the generator.
  The admissible values are {\tt (Full, Recurrent, Orbit, PrimePower)}.
  \begin {description}
  \item [{\tt Full} :] The complete lattice, generated by all possible
     initial states, will be analyzed.
  \item [{\tt Recurrent} :] If the (combined) generator has transient
     states, then the lattice analyzed will be the sublattice generated by
     the set of recurrent states.
  \item[{\tt Orbit} :] The grid generated by the (forward) orbit of 
     a state of the (combined) generator is analyzed. This state
     is specified as follows. 
     On the following $J$ lines, the initial state for each
     component must be given. This is an integer vector with a number 
     of components equal to the order of the component.
  \item [{\tt PrimePower} :] In the case where some component is an MLCG
     whose modulus is a power of a prime $p$, then the states visited over
     a single orbit (subcycle) of that component generate a sublattice
     (when $a \equiv 1 \pmod{p}$) or belong to the union of
     $p-1$ sublattices (otherwise).  If {\em LatticeType\/} takes this value,
     if a component is an MLCG ($k=1$), 
     and if the modulus of that MLCG is given
     in the data file in the form (b): $(x\ y\ z)$ with $z=0$ and $x$ prime, 
     then what is analyzed
     is one of those sublattices.  This is done by dividing the modulus
     by the appropriate power of $p$, as described in \cite{rLEC97c}.
     For example, if $p=2$ and $a \mod 8 = 5$, then the modulus is divided
     by 4 as in \cite{rFIS90a,rKNU81a}.
  \end {description}
%
\item [{\em LaGroupSizes, ~Spacing\/} :]
  These data fields are positive integers, used to introduce lacunary indices.
  If the respective values are $s$ and $d$, then we will analyze the
  lattice structure of vectors of the form
  $(u_{i+1}, \dots, u_{i+s}, u_{i+d+1},\dots, u_{i+d+s}, u_{i+2d+1},\dots,
  u_{i+2d+s}, \dots)$, formed by groups of $s$ successive values,
  taken $d$ values apart.
  To analyze vectors of successive values (as usual), take $s=d=1$
  or $s$ larger or equal to {\em MaxDim}.
  To analyze lacunary indices that are not evenly spaced,
  put $s = -t$ where {\em $t=$MaxDim\/} and then, on the $t$ lines that
  follow, give the $t$ lacunary indices $i_1,\dots,i_t$, which are to
  be interpreted as in Section 1.4.
%
\item [{\em VerifyBB} :]
  Indicates whether or not the results of the branch-and-bound procedures
  will be verified with formal error bounds on the numerical computational
  errors.  The verification (if any) is performed using either
  {\tt REDBAS2} or {\tt REDBAS3}.
  When such a verification is not performed, the results are not
  formally guaranteed, but computations are performed an order of
  magnitude faster.  The possible values are:
 \begin {description}
 \item [{\tt VerifyAll2} :] Performs the verification for all the generators
   examined, using {\tt REDBAS2}.
 \item [{\tt VerifyRetained2} :] Performs the verification only for the
   {\em NbGen\/} ``best'' generators retained, using {\tt REDBAS2}.
 \item [{\tt VerifyAll3} :] Performs the verification for all the generators
   examined, using {\tt REDBAS3}.
 \item [{\tt VerifyRetained3} :] Performs the verification only for the
   {\em NbGen\/} ``best'' generators retained, using {\tt REDBAS3}.
 \item [{\tt NoVerify} :] Performs no verification at all.
 \end {description}
%
\item [{\em MaxNodesBB} :]
  An integer giving the maximum number of nodes to be examined in any
  given branch-and-bound procedure when computing $d_t$ or $q_t$.
  When that value is exceeded, the branch-and-bound is stopped and the
  generator is rejected.  The number of generators rejected for that
  reason is given in the results.
  A small value of {\em MaxNodesBB} will make the program run faster
  (sometimes much faster), permitting to examine more generators,
  but will increase the chances of rejecting good generators.
%
\item [{\em TimeLimit} :]
  A real number giving the maximal CPU time, in hours, given to the program
  for performing its search.
  When this time-limit expires, the partial results are printed, with a message.
  For a random search, the number of subregions printed includes
  the last region searched (whose search may not be finished).
%
\item [$S1,\, S2$ :]
  Those integers are the seeds of the generator used during the
  search, if necessary; e.g., to generate the $\alpha_i$'s.
  Before starting the search, the program calls
  {\tt SetSeed($S1, S2$)} form the module {\tt SUPRAND} of SENTIERS
  and the module {\tt RAND} of SIMOD.
  To perform a different random search in a region already studied,
  just change these seeds.
  One must have $1\le S1\le 2147483562$ and $1\le S2\le 2147483398$.
%
\item [{\em ResultForm} :]
  Selects in which form the results will be given.
  The possible values are {\tt (Terminal, RES, GEN, RESGEN)}.
  Lowercases are also allowed.
  \begin {description}
   \item [{\tt Terminal}] indicates that the results will appear only on the
     terminal screen.
   \item [{\tt RES}] says that the results will be in a file with the same
     name as the data file, but with extension ``{\tt .res}''.
   \item [{\tt GEN}] says that the retained generators will be listed in
     a file with the same name as the data file, with extension ``{\tt .gen}''.
     This file can then be taken as input to the same program, for
     example to perform a second pass with a more stringent criterion
     or to compute higher dimensional lattice ``measures'' for the
     retained generators.
   \item [{\tt RESGEN}] means that the results will be in a ``{\tt .res}''
     file {\em and} the retained generators will be listed in a ``{\tt .gen}''
     file, as with {\tt RES} and {\tt GEN} above.
  \end {description}
\end {description}


\defmodule {lats}

This is the same program as {\tt latl}, except that multipliers and
bases are represented using the type {\tt SuperInteger}.
(Modules {\tt MULTSI} and {\tt BASISSI} are used.)
Therefore, it can deal with larger values of $m$ and $a_i$'s.

\bigskip\bigskip
%%%%%%%%%%%%%%%%%%%%%%%%%%%%%%%%%
\defmodule {latsr}

Same as {\tt latl}, except that the multipliers
are represented as {\tt SuperInteger}'s.
(Modules {\tt MULTSI} and {\tt BASISLR} are used.)

\bigskip\bigskip
%%%%%%%%%%%%%%%%%%%%%%%%%%%%%%%%%
\defmodule {seeks}

Same as {\tt seekl}, except that the multipliers and
bases and represented as {\tt SuperInteger}'s.
(Modules {\tt MULTSI} and {\tt BASISSI} are used.)

\bigskip\bigskip
%%%%%%%%%%%%%%%%%%%%%%%%%%%%%%%%%
\defmodule {seeksr}

Same program as {\tt seekl}, except that multipliers
are represented as {\tt SuperInteger}'s.
(Modules {\tt MULTSI} and {\tt BASISLR} are used.)
It can deal with larger values of $m$ and $a_i$'s than {\tt seekl},
but the computations are less precise than with {\tt seeks}.

\bigskip\bigskip
\input streams

  % This includes  lats, latsr, seeks, seeksr, streams.
%\input streams
\defmodule {menumrg}

This is an interactive program to build, manipulate, compare, and
reduce lattice bases associated with LCGs or MRGs.
It uses the ``{\tt SI SI}'' representation.
The main purpose is to examine interactively what happens to a
basis when different operations are performed on its vectors.
This could be useful for ``debugging'' or to deal with special classes
of bases for which the standardized procedures may not work very well.
A variant of this module has been used, for instance, to obtain the
numerical results given in \cite{rCOU94a,rTEZ93a}.
In those cases, the shortest vectors and reduced bases were
computed in (partly) ``manual'' mode.

Two bases {\tt B1} and {\tt B2} and one vector {\tt Ve}
are created initially (automatically) and available to the user.
All along, {\tt B1} and {\tt B2} will be bases of the same lattice.
The program performs different operations, classified according to which
of those three objects they involve:
(a) the operations which affect only {\tt B1};
(b) those affecting {\tt B1} and {\tt B2};
(c) and those involving {\tt B1} and {\tt Ve}.
All those operations are available from a menu (and some submenus), as
shown in Figure~\ref{mainmenu}.
We now briefly explain the different menu options.

%%%%%%%%%%%%%%%%%%%%%%%%%%%%%%%%%%%%%%%%%%%%%%%%%%%%%%%
\setbox1=\vbox {\hsize = 6.0in
\code

  Operations on B1 :
    1. Building MRG basis
    2. Reading/Writing
    3. Duality verification
    4. Dualizing
    5. Square lengths
    6. Sorting
    7. Pre-reduction
    8. Shortest vector/Minkowski reduction
    9. Shortest/Minkowski verification

  Operations involving B1 and B2 :
    10. Copy B1 in B2
    11. Permute B1 and B2
    12. Equivalence of B1 and B2
    13. Relative coordinates
    14. Lexicographic order


  Operations involving B1 and Ve :
    15. Vector construction
    16. Relative coordinates
    17. Insertion of Ve in B1

\endcode
}

\begin{figure}[htb] \centering \boxit{\box1}
\caption { Main menu for the program {\tt menumrg}.}
\label {mainmenu}
\end{figure}


\begin {description}
%
\item [{\bf Operations on B1 :}]
\item [1. {\em Building MRG basis\/} :]
   Asks the user for a modulus, order, multiplier(s), and dimensions,
   then constructs the appropriate lattice basis in {\tt B1}.
\item [2. {\em Reading/Writing\/} :]
   Read/write the basis {\tt B1} from/to a file or the terminal.
\item [3. {\em Duality verification\/} :]
   Verifies whether the dual basis is ``correct'' by computing the
   appropriate scalar products.
\item [4. {\em Dualizing\/} :]
   Interchange the primal/dual bases.
\item [5. {\em Square lengths\/} :]
   Compute the squared vector lengths.
   The values can be displayed either in full (large integers)
   or in scientific notation.
\item [6. {\em Sorting\/} :]
   Sorts the vectors of {\tt B1} by increasing length.
\item [7. {\em Pre-reduction\/} :]
   Performs ``pairwise'' prereductions $T$ and $T^*$ of \cite{rDIE75a}
   in various combinations or orders.
   Can also ask for a given number of pre-reductions.
\item [8. {\em Shortest vector/Minkowski reduction\/} :]
   Computes the shortest vector in the lattice, or a MRLB,
   using the module {\tt REDBAS} (does not compute error bounds).
\item [9. {\em Shortest/Minkowski verification\/} :]
   Verifies whether the first vector of {\tt B1} is really a
   shortest vector, or whether the current basis is really a MRLB,
   by performing the branch-and-bound with error bounds
   (using either {\tt REDBAS2} or {\tt REDBAS3}).
   An option is also available to compute the squared distance between
   each vector of {\tt B1} and the subspace generated by all the other vectors,
   and to select the vectors for which that distance is smaller than
   the length of the (current) shortest vector in {\tt B1}.
%
\item [{\bf Operations involving B1 and B2 :}]
\item [10. {\em Copy B1 onto B2\/} :]
   Copies the basis {\tt B1} onto {\tt B2}.
\item [11. {\em Permute B1 and B2\/} :]
   Permutes {\tt B1} and {\tt B2}.
\item [12. {\em Equivalence of B1 and B2\/} :]
   Tells whether or not {\tt B1} and {\tt B2} generate the same lattice.
   Assumes that the dual basis is ``updated''.
\item [13. {\em Relative coordinates\/} :]
   Gives the relative coordinates of a vector of {\tt B1} relative to
   the basis {\tt B2}.
\item [14. {\em Lexicographic order\/} :]
   Displays the squared lengths of the vectors of {\tt B1} and {\tt B2}
   in lexicographic order.
   Assumes that the two bases are already sorted.

%
\item [{\bf Operations involving B1 and Ve :}]
\item [15. {\em Vector construction\/} :]
   Permits one to either initialize {\tt Ve} to zero; 
   add to {\tt Ve} a multiple of a vector of {\tt B1};
   increase the dimension of {\tt Ve} by one (the new coordinate is zero);
   shift the coordinates of {\tt Ve} to the right by one (provided that the
   last coordinate is zero); or perform a linear combination of iterates
   of the latter shifting operations (as in \cite{rCOU94a}).
\item [16. {\em Relative coordinates\/} :]
   Computes the coordinates of {\tt Ve} relative to {\tt B1}.
\item [17. {\em Insertion of Ve in B1\/} :]
   Inserts {\tt Ve} into the basis {\tt B1}, provided that {\tt Ve}
   belongs to the lattice generated by {\tt B1} and is primitive.
\end {description}

%\defmodule{palpha}

\baselineskip=14pt

Le $P_\alpha(s)$ est une mesure de la qualit\'e d'int\'egration d'une
r\`egle de treillis.  D'une mani\`ere g\'en\'erale, 
une r\`egle de treillis permet d'approximer
l'int\'egrale de $f$, une fonction, sur le cube unitaire de dimension $s$.
$f$ doit \^etre lisse et p\'eriodique par rapport \`a $x \in$ {\bf R}.

La formule g\'en\'erale du $P_\alpha(s)$ pour $\alpha$ entier et pair est 
\eq
  P_\alpha (s) = -1 + 1/m \sum_{i=0}^{m-2} \prod_{j=0}^{s-1} \Big[
                     1- \frac{(-1)^{\alpha /2}(2\pi)^{\alpha}}{\alpha !}B_\alpha(x_{i+j}) \Big] 
										\label{eq:Palpha}
\endeq
o\`u les $B_\alpha(x)$ sont les polyn\^omes de Bernoulli. $x_{i+j}$ est le $i+j^e$ nombre
g\'en\'er\'e et se situe entre 0 et 1.
Les premiers polyn\^omes de Bernoulli sont \\\\ 
\hspace*{18mm} \baselineskip=18pt
$B_0(x)=1$,\hspace*{9mm} $B_1(x)=x-1/2$,\hspace*{9mm}
$B_2(x)=x^3-x^2+1/6$,  \\ \hspace*{18mm}
$B_3(x)=x^3-(3/2)x^2+(1/2)x$,\hspace*{5mm} $B_4(x)=x^4-2x^2+x-1/30$

\noindent
Les $B_\alpha(x)$ se calculent en utilisant l'extension de la s\'erie suivante:
$$\frac{te^{xt}}{e^t-1} = \sum_{s=0}^{\infty}B_s \frac{t^s}{s!}$$
On peut ajouter un poids \`a chacune des dimensions $s$, pour des $\alpha$ pairs.
La formule du
$P_\alpha(s)$ devient: 
$$P_{2\alpha}(s) = \beta_0 \Big\{ -1 + 1/m \sum_{i=0}^{m-2} \prod_{j=0}^{s-1} \Big[
                     1- \frac{\big( -(2\pi\beta_{j+1})^2\big) ^{\alpha}}{2\alpha !}B_{2\alpha}(x_{i+j}) \Big] \Big\} 
$$

Le $B_\alpha(s)$, \`a ne pas confondre avec le $B_\alpha(x)$, est une mesure
indiquant qu'il existe au moins un $P_\alpha(s) \le B_\alpha(s)$.  D'une mani\`ere
g\'en\'erale, on peut le voir comme une borne sup\'erieure valide en pr\'esence
des conditions suivantes: $m> e^{\alpha s/(\alpha -1)}$ et premier, $s\ge 2$ et
$\alpha > 1$.  Pour plus de d\'etails consulter le th\'eor\`eme 4.4, p.83 du
livre de Sloan et Joe. 
Il se calcule ainsi:
\begin{eqnarray*}
B_\alpha(s) = {\Big( \frac{e}{s} \Big)}^{\alpha s} \frac{(2\ln m + s)^{\alpha s}}{m^\alpha}
\end{eqnarray*} 

\large
{\bf Le module PALPHA}
\normalsize

Module servant de support aux programmes PAl, BAl et Coeff. Ils permettent
respectivement de trouver le $P_\alpha$ ou  le $B_\alpha$ d'un g\'en\'erateur
ou enfin de trouver des bons coefficients pour un g\'en\'erateur en
utilisant le $P_\alpha$ comme discriminant.  

\large
{\bf palpha}
\normalsize

Ce programme permet de calculer le $P_\alpha$ pour un 
g\'en\'erateur et un $\alpha$ donn\'e. 
Tout d\'ependant du
$\alpha$
la fonction de calcul changera. Les donn\'ees suivantes doivent \^etre
plac\'ees dans un fichier avec l'extension {\tt".dat"} selon le format de
la figure~1.  Les r\'esultats seront plac\'es dans un fichier  avec l'extension {\tt".res"}.

\large
{\bf balpha}
\normalsize

Permet de calculer le $B_\alpha$ d'un g\'en\'erateur.  Pour l'instant
seulement le cas avec $\alpha = 2$ est impl\'ement\'e.  Le format du
fichier est le m\^eme que pour le programme PAl. Les donn\'ees devront
\^etre plac\'ees dans un fichier avec l'extension {\tt".dat"} et les
r\'esultats se trouveront dans le fichier avec le m\^eme pr\'efixe,
cette fois avec l'extension {\tt".res"}.

\noindent
{\bf Le format du fichier d'entr\'ee pour PAl ou BAl}

\begin{center}
\it{
\begin{tabbing}
\hspace{2cm}\= LCG \hspace{3cm} \= SorteGenerateur \\
\>PAl ou BAl \>   Type de calcul \\
\>251     \>      m  modulo\\
\>true    \>      modulo premier\\
\>33      \>      a  multiplicateur\\
\>2       \>      mindim\\
\>4      \>      maxdim\\
\>2      \>      saut \\ 
\>2       \>      alpha\\
\>false   \>      VerifBooleen\\
\>1       \>      seed\\
\>1.0 2.0 3.0 1.0 1.0 \> Beta(0),...,Beta(MaxDim)\\
\end{tabbing}
}
\rm {Figure 1 : Format du fichier pour PAl ou BAl}
\end{center}

%%\begin{figure}[h]
%%\caption{Format du fichier pour PAl ou BAl}
%%\end{figure}

\noindent
{\it Type de calcul} : Permet de d\'eterminer le type de calcul
\`a effectuer.

\noindent
{\it Modulo} : le modulo du g\'en\'erateur en base 10, soit un nombre premier ou
une puissance de 2.  Ce nombre doit \^etre inf\'erieur \`a 3.03E09 ou $2^{31}$.

\noindent
{\it Modulo premier} : Bool\'een indiquant si le modulo utilis\'e est un 
nombre premier.

\noindent
{\it Multiplicateur} : le multiplicateur du g\'en\'erateur en base 10. Il doit \^etre
inf\'erieur au modulo.  Ou un intervalle dans le cas de la recherche de coefficient.
La recherche permettra de trouver le meilleur coefficient dans cet intervalle en fonction du
$P_\alpha$.

\noindent
{\it Mindim} : Entier  donnant la borne inf\'erieur sur le s (la dimension) 
dans le calcul \`a effectuer.  

\noindent
{\it Maxdim} : Entier  donnant la borne sup\'erieur sur le s (la dimension)
dans le calcul \`a effectuer.

\noindent
{\it Saut} : Entier d\'eterminant le saut \`a effectuer entre les dimensions.

\noindent
{\it alpha} : Entier pair, en g\'en\'eral de 2 \`a 10.

\noindent
{\it VerifBooleen} : Si {\tt true} alors on v\'erifira si la p\'eriode est maximale,
mettre \`a {\tt false} si aucune v\'erification n'est n\'ecessaire.  Dans le cas de 
la recherche de coefficient, si {\tt true} alors seulement les g\'en\'erateurs
\`a p\'eriode maximale seront consid\'er\'es, si {\tt false} alors le
$P_\alpha$ de tous les g\'en\'erateurs sera \'evalu\'e.  Si par contre le
modulo n'est pas un nombre premier alors peu importe la valeur entr\'ee \`a
{\it VerifBooleen},
nous allons effectuer une recherche sur l'ensemble des coefficients.


\noindent
{\it Seed} : Entier inf\'erieur au modulo qui servira de germe dans le g\'en\'erateur.  

\noindent
{\it Beta} : Nombres r\'eels permettant de donner un poids \`a chacune des dimensions.
Il faut entrer dans le fichier une valeur pour $\beta (0),...,\beta (MaxDim)$. 
  % Now a special case of latl and lats.
%\defmodule{coeff}

\baselineskip=14pt
Ce programme permet de trouver des bons coefficients pour un
g\'en\'erateur dont on a donn\'e le modulo et le type. Cette
recherche se fait en fonction des $P_\alpha$. Il faut donc 
\'egalement pr\'eciser un $\alpha$. 

\noindent
{\bf Le format du fichier d'entr\'ee pour Coeff}

\begin{center}
\it{
\begin{tabbing}
\hspace{2cm}\= LCG \hspace{3cm} \= SorteGenerateur \\
\>Coeff \>   Type de calcul \\
\>251     \>      m  modulo\\
\>true    \>      modulo premier\\  
\>3 250   \>      intervalle pour le multiplicateur a\\
\>2       \>      mindim\\
\>4      \>      maxdim\\
\>2      \>      saut \\ 
\>2       \>      alpha\\
\>false   \>      VerifBooleen\\
\>1       \>      seed\\
\>1.0 2.0 3.0 1.0 1.0 \> Beta(0),...,Beta(MaxDim)\\
\end{tabbing}
}
\rm {Figure 1 : Format du fichier pour Coeff}
\end{center}

\noindent
Tous les champs sont identiques au cas du PAl ou BAl sauf pour le
cinqui\`eme qui indique plut\^ot l'intervalle \`a l'int\'erieur 
duquel chercher le meilleur coefficient.
   % Now a special case of seekl and seeks.
%
%                   Intermediate-level modules.
\thispagestyle {nomark}
\addcontentsline{toc}{section}{\protect\numberline{B.}{
   Intermediate-Level Modules }}
\null\vfill\vfill
\centerline {\large\bf APPENDIX B}
\vfill
\centerline {\large\bf Intermediate-Level Modules}
\vfill\vfill
%\eject\null\vfill         % Page blanche.

\defmodule{LATMRG}

The files {\tt LATMRG.h} and  {\tt LATMRG.c} exist uniquely to permit one to 
use LatMRG directly from a {\bf C} program.   For this, one must include the directive

\noindent {\tt \#include "LATMRG.h"}

\noindent at the beginning of the file and then add

\noindent {\tt LATMRG\_BEGIN (argc, argv);}

\noindent as a first instruction in the fonction 

\noindent {\tt int main(int argc, char **argv);}

To use functions from {\tt mylib}, {\tt simod}, and {\tt sentiers}, 
one must include the header files

\noindent {\tt \#include "MYLIB.h"}\hfill\break
\noindent {\tt \#include "SIMOD.h"}\hfill\break
\noindent {\tt \#include "SENTIERS.h"}

One can consult the appropriate {\tt *.h} header files in {\tt latmrg/xds} 
to see the prototypes for the external {\bf C} functions.


\bigskip\hrule
\code
\hide
<* +M2EXTENSIONS *>
<* +NOHEADER *>
<* +CSTDLIB *>
\endhide
DEFINITION MODULE ["C"] LATMRG;

PROCEDURE LATMRG_BEGIN();

END LATMRG.
\endcode

\include {PALPHA}
\defmodule {LATIO}

This module contains several tools for input/output,
for performing generator searches, for managing and storing lists
of generators, and the like.
It contains the actual implementations of the
high-level ``seek'' and ``lat'' programs.
The procedures here use most of the other 
lower-level and intermediate-level modules.
All the {\tt Read...} procedures read the data in the format
explained in the documentation of the {\tt seekl} program.

The searches for lists of (good or bad) generators are 
performed by categories, as explained in
the documentation of the {\tt seekl} program.
The set of dimensions for which the lattice structure is examined
is partitioned into a certain number of intervals (one or more),
also called {\em categories\/}.
We give bounds on the accepted values of the figure of merit
for each category and, for any given category $c$,
we consider only those generators that satisfy the bounds 
for all the categories less or equal to $c$.
Then we retain a list of generators with the best or worst figures of
merit for each category.


\bigskip\hrule
\code

DEFINITION MODULE LATIO;

FROM SUPINT   IMPORT SuperInteger;
FROM SUPFACT  IMPORT Factors;
FROM CONFIG   IMPORT BDim, OutputType, MaxJ, MaxCat, Indexj, MaxNbGen;
FROM MULT     IMPORT MScal, MVect;
FROM NORM     IMPORT CriterionType; 
FROM REDBAS   IMPORT VerifyType; 
IMPORT BASISLR;


TYPE
   GenerType   = (MRG, MWC);
\endcode
 \tabb Type of generator that can be considered. \endtabb
\code

   CondType    = (NoCond, AppFact, MaxBits);
\endcode
 \tabb Type of condition on the multipliers. \endtabb
\code

   VerifType   = (VerifyAll2, VerifyRetained2, 
                  VerifyAll3, VerifyRetained3, NoVerify);
\endcode
 \tabb Type of verification that can be performed on the results
    (i.e., type of computational method). \endtabb
\code

   Poly        = POINTER TO InfoPoly;
   InfoPoly    = RECORD
      a        : MVect;
      NextPoly : Poly
      END;
\endcode
 \tabb  Elements to store a linked list of polynomials.  \endtabb
\code

   Resultats   = POINTER TO InfoRes;
   InfoRes     = RECORD
      RD2, RS2, RQ2 : BASISLR.BVect;
      DimD2, DimQ2  : BDim
      END;
\endcode
 \tabb  Type of record used as a field of the record type {\tt InfoGen},
    to store detailed results for a given generator. 
    It contains the {\em squared\/} values of the $d_t$, $S_t$, and $q_t$,
    and the maximal dimensions {\tt DimD2} and {\tt DimQ2} 
    for which $d_t$ and $q_t$, respectively, have been computed.
 \endtabb
\code

   Generateur = POINTER TO InfoGen;
   InfoGen    = RECORD
      Maj      : ARRAY Indexj OF MVect; 
      a        : MVect;
      DimMerit : BDim; 
      Merit    : LONGREAL;
      Res      : Resultats;
      NextGen  : Generateur
      END;
\endcode
 \tabb  Type of record used to store a linked list of generators.
   Used to store the generators that are kept in each category.
   The coefficients $a_{j,i}$ of the components are in {\tt Maj}
   whereas those of the combined generator are in {\tt a}.
   {\tt DimMerit} is the dimension where the minimum is attained for
   the figure of merit (for the given category) and {\tt Merit} gives
   the value of this figure of merit.
   {\tt Res} is a pointer to the more detailed results and
   {\tt NextGen} is a pointer to the next generator in the list.
 \endtabb
\code

   ToZone = POINTER TO Zone; 
   Zone   = RECORD
      No         : INTEGER; 
      Binf, Bsup : SuperInteger;
      p          : LONGREAL;
      PetiteDiff : BOOLEAN;
      BsupMsH    : SuperInteger;
      Mq, MBsup  : MScal;
      NextZone   : ToZone;
      END;
\endcode
 \tabb  Type of record used to store the research zones.
   {\tt No} is the number of the zone, which is in $\{0,1,2,3\}$.
   {\tt Binf} and {\tt Bsup} are the lower and upper bounds
   defining the zone.  
   {\tt P} is the fraction of the acceptable
   values of $a$ lying in this zone (this is used for the random searches).
   {\tt PetiteDiff} is {\tt TRUE} iff {\tt Bsup - Binf $\le$ H} (or {\tt Hk}),
   and when it is {\tt TRUE}, {\tt BsupMsH} is equal to 
   {\tt Bsup ${}+1-{}$H} (or {\tt Hk}). This is also used for random searches.
   The variables {\tt Mq} and {\tt MBsup} are used in {\tt ExamAllZones}.
  \endtabb
\code

   ToRegion = POINTER TO Region;
   Region   = RECORD
      No         : INTEGER;
      Binf, Bsup : MScal;
      Mq         : MScal
      END;
\endcode
 \tabb  Type of record used to store the regions that are to be searched.
   {\tt No} is the zone number where this region is (in $\{0,1,2,3\}$).
   {\tt Binf} and {\tt Bsup} are the lower and upper bounds
   defining the region.  
   The variable {\tt Mq} is  used in {\tt ExamRegion}.
  \endtabb
\code

PROCEDURE TimeOver () : BOOLEAN;
\endcode
 \tab 
   Returns {\tt TRUE} if and only if the cpu time limit has been reached.
 \endtab
\code

PROCEDURE CreateGenerateur (VAR G : Generateur);
\endcode
 \tab 
  Creates a record for a generator {\tt G} and initializes its fields.
 \endtab
\code

PROCEDURE CreateZone (VAR Z : ToZone);
\endcode
 \tab 
  Creates a record for a zone {\tt Z} and initializes its  fields.
 \endtab
\code

PROCEDURE CreateRegion (VAR R : ToRegion);
\endcode
 \tab 
  Creates a record for a region {\tt R} and initializes its  fields.
 \endtab
\code

PROCEDURE LireFact (VAR L : Factors;   NomFch : ARRAY OF CHAR);
\endcode
 \tab
  Reads a list of factors in the file {\tt NomFich} and places them
  in the list {\tt L}.  Assumes that the factors in the file are
  written according to the format used by {\tt Ecrirefact}.
 \endtab
\code

PROCEDURE EcrireFact (L : Factors;   NomFch : ARRAY OF CHAR);
\endcode
 \tab
  Writes the factors of the list {\tt L} to the file {\tt NomFich}.
  It this file exists, it is rewritten.
  The factors can be read back later by {\tt LireFact}.
 \endtab
\code

PROCEDURE Readm (j : Indexj);
\endcode
 \tab
  Reads the modulo $m_j$ for the component $j$.
  Initializes some constants that depend on it.
 \endtab
\code

PROCEDURE ReadOnea (aa : SuperInteger);
\endcode
 \tab
  Reads one multiplier $a_{j,i}$ or one bound 
  $b_{j,i}$ or $c_{j,i}$ in the format described in {\tt seekl}.
 \endtab
\code

PROCEDURE ReadReadGen;
\endcode
 \tab
  Reads a boolean indicating if the generators to be tested should be
  taken from a `{\tt .gen}' file, and if true, \
  also reads the name of the file, without the {\tt .gen} extension.
 \endtab
\code

PROCEDURE ReadTypeGen (j : Indexj);
\endcode
 \tab
  Reads the type of generator to be considered, MRG or MWC,
  for component $j$.
 \endtab
\code

PROCEDURE ReadOrder (j : Indexj);
\endcode
 \tab
  Reads the order $k_j$ of the recurrence  for component $j$.
 \endtab
\code

PROCEDURE ReadPerMax (j : Indexj);
\endcode
 \tab
  Reads a boolean indicating if component $j$ should have 
  maximal period.  The muliplier $m_j$ must have been read in
  the `3 numbers' format by {\tt readmj}, otherwise this boolean
  is reset to {\tt FALSE}.
 \endtab
\code

PROCEDURE ReadImplemCond (j : Indexj);
\endcode
 \tab
  Reads the condition that we impose on the coefficients $a_{j,i}$
  for the component $j$.  
 \endtab
\code

PROCEDURE ReadFactmr (j : Indexj);
\endcode
 \tab
  Reads what to do about the factorizations of $m_j-1$ and 
  $r_j = (m_j^{k_j}-1)/(m_j-1)$,  for the component $j$.
 \endtab
\code

PROCEDURE Reada (j : Indexj);
\endcode
 \tab
  Reads the $k_j$ multipliers $a_{j,1},\dots,a_{j,k_j}$ for component $j$.
 \endtab
\code

PROCEDURE Readbc (j : Indexj);
\endcode
 \tab
  Reads the $2k_j$ bounds $b_{j,1},c_{j,1}\dots,b_{j,k_j},c_{j,k_j}$ 
  on the multipliers for component $j$.
 \endtab
\code

PROCEDURE ReadSearchMethod (j : Indexj);
\endcode
 \tab
  Reads the method of search, exhaustive or random, for the polynomials
  for component $j$.  If the method is {\tt Random}, the it reads also
  the number of regions and their sizes.
 \endtab
\code

PROCEDURE ReadMRGComponent (j : Indexj);
\endcode
 \tab
  Reads all the parameters related to the component $j$,
  by calling the appropriate `{\tt Read...}' procedures above.
  Also creates the required data structures for this component.
 \endtab
\code

PROCEDURE ReadMinMaxDim;
\endcode
 \tab
  Reads the minimal and maximal numbers of dimensions in which the
  test is to be applied, the number of categories,
  and the boundaries of these categories.
  The format is slightly different for the `seek' and `lat'
  programs.
 \endtab
\code

PROCEDURE ReadMinMaxMerit;
\endcode
 \tab
  Reads the minimum and maximum acceptable values for the figure 
  of merit for each category of dimensions.
 \endtab
\code

PROCEDURE ReadNbGen;
\endcode
 \tab
  Reads the number of generators that we want to retain in each
  category.
 \endtab
\code

PROCEDURE ReadCriterion;
\endcode
 \tab
  Reads the selection criterion (figure of merit).
 \endtab
\code

PROCEDURE ReadLatInfo;
\endcode
 \tab
  Reads what lattice information we want to be printed in the results.
 \endtab
\code

PROCEDURE ReadLatType;
\endcode
 \tab
  Reads which type of lattice we want to analyze:
  That generated by the recurrent states only, or that generated by
  a single orbit, or that generated by the longest cycle of a 
  generator with prime power modulus, or the full lattice generated
  by the union of all the orbits.
  In the case of {\tt Orbit}, this procedure also reads the initial
  state and initializes some additional data structures.
 \endtab
\code

PROCEDURE ReadLaIndices;
\endcode
 \tab
  Reads the indices (lacunary or successive) that are to be considered.
 \endtab
\code

PROCEDURE ReadVerifyBB;
\endcode
 \tab
  Reads the verification level that we want to use.
 \endtab
\code

PROCEDURE ReadMaxNodesBB;
\endcode
 \tab
  Reads the maximum number of nodes that will be allowed in
  the branch-and-bound tree.  When that number is exceeded,
  the branch-and-bound procedure will stop.
 \endtab
\code

PROCEDURE ReadTimeLimit;
\endcode
 \tab
  Reads the cpu time limit that is allowed for a search.
  When that limit is exceeded, the search procedure stops
  and reports what it has found so far.
 \endtab
\code

PROCEDURE ReadSeeds;
\endcode
 \tab
  Reads the seeds that are used for the random number generator
  of SENTIERS.  This generator is used for the random searches.
 \endtab
\code

PROCEDURE ReadResultForm;
\endcode
 \tab
  Reads the form of the output (on the terminal, in a {\tt .gen} file,
  in a {\tt .res} file, in a {\tt .tex} file).
 \endtab
\code

PROCEDURE Computem;
\endcode
 \tab  Computes the order and the modulus $m$ for the 
   combined generator, given those of the components.
   Computes also the period lengths of the components and of
   the combination.
 \endtab
\code

PROCEDURE Computea;
\endcode
 \tab 
   Computes the vector of multipliers $a$ for the combined generators,
   given the parameters of the components.
   Assumes that the moduli of the components are all prime.
   Initializes the vectors {\tt Maj} and {\tt a}. 
 \endtab
\code

PROCEDURE SaisieDonnees;
\endcode
 \tab
  Calls the different low-level procedures to read the data
  from the data file, in the format appropriate for the programs
  {\tt latl} and {\tt seekl}.
  Also creates some of the data structures (e.g., {\tt SuperInteger}'s)
  that are required.
 \endtab
\code

PROCEDURE ReadGenFile;
\endcode
 \tab
  Reads a set of generators from the file `{\tt xxx.gen}',
  assuming that the data has been read from `{\tt xxx.dat}'. 
 \endtab
\code

PROCEDURE PerMaxPowPrime (j : Indexj) : BOOLEAN;
\endcode
 \tab
   Assumes that $m_j$ is a power of a prime $p$, that $k_j = 1$, 
   and that the recurrence is homogeneous.   
   Returns {\tt TRUE} iff the maximal period conditions are 
   satisfied for that case for component $j$.
 \endtab
\code

PROCEDURE Reducemj (VAR a : MVect);
\endcode
 \tab
  Reduces $m_j$ and reinitializes the module {\tt MRG} accordingly
  in the case where the modulus is a power of a prime and where
  {\tt LatType = PrimePower}.  Called by {\tt testGen}.
 \endtab
\code

PROCEDURE InitZones (j, i : INTEGER);
\endcode
 \tab
  Initializes the research zones for the multiplier $a_i$ of the
  component $j$.
  In the case of a random search, this procedure also creates the
  region for $(j,i)$.    Called by {\tt SetupGenSeek}.
 \endtab
\code

PROCEDURE ChoisirBornes (j : Indexj);
\endcode
  \tab  This procedure is called by {\tt SeekGen} in the case of a
    random search.  It chooses a region at random and 
    initializes its boundaries.  This region will then be searched
    completely.  
  \endtab
\code

PROCEDURE TestGen
   ( VAR a    : MVect;
     T        : CriterionType;
     VerifyBB : VerifyType;
     Out      : OutputType
   );
\endcode
 \tab
   Tests the (combined) generator with multiplier (vector) {\tt a}.
 \endtab
\code

PROCEDURE ConserverGen (C : INTEGER;  tMerit2 : BDim);
\endcode
 \tab
   Inserts the most recently tested generator into the list
   for the categoty {\tt Cat}.
   The value {\tt tMerit2} indicates the dimension where the
   worst-case is attained for the figure of merit.
   This procedure is called by {\tt VerifyCategories}.
 \endtab
\code

PROCEDURE VerifyCategories (VAR Me2 : BASISLR.BVect);
\endcode
 \tab
  Verifies if the most recently tested generator can be included
  in any of the lists corresponding to the different categories.
  If so, includes it in the appropriate lists. 
  The vector {\tt Me2} is either {\tt Q2} or {\tt S2}, depending
  on the figure of merit that is used.
 \endtab
\code

PROCEDURE ExamAllZones (j, i : INTEGER);
\endcode
 \tab
  This procedure is used in the ``exhaustive search'' case. 
  When $J > 1$, it collects primitive polynomials for the given $j$.
  When $J = 1$, it tests all generators in the region.
  The initial call should be with $i$ equal to the degree of the
  polynomials for component $j$.
  Then, the procedure calls itself recursively for smaller $i$.
 \endtab
\code

PROCEDURE ExamRegion (j, i : INTEGER);
\endcode
 \tab
  Similar to {\tt ExamAllZones}, but used for the ``Random search'' case.
 \endtab
\code

PROCEDURE ExamCombPoly (jj : Indexj);
\endcode
 \tab
  Used only when $J > 1$.  Examines all combinations of the primitive
  polynomials found by {\tt ExamRegion} or {\tt ExamAllZones}.
 \endtab 
\code

PROCEDURE ExamCurrenta;
\endcode
 \tab
  Tests the generator that corresponds to the current value of the
  multiplier (vector) {\tt a}, and adds it to the list(s) if it is competitive.
  Called only by {\tt TestGenFromFile}.
 \endtab
\code

PROCEDURE TestGenFromFile;
\endcode
 \tab
  Examines all the generators from a {\tt .gen} file,
  which is assumed to be the current input file.
 \endtab
\code

PROCEDURE CalculNbreGen (j : Indexj);
\endcode
 \tab
  Computes the total number of simple generators (i.e., polynomials)
  examined, for the component {\tt j}, and puts the result in
  {\tt TotGen [j]}.
 \endtab
\code

PROCEDURE CalculerResteInfo;
\endcode
 \tab 
  After the search has been completed, computes all the remaining required
  information on the retained generators.
  Computes the $d_t$ or the $q_t$ when required and not already done.
  Redo the computations with error bounds if requested.
 \endtab
\code

PROCEDURE PrintErrBounds (T : INTEGER);
\endcode
 \tab
  Computes and prints the bounds on the Euclidean distances between
  the points produced by a L'Ecuyer-style combined MRG and the 
  corresponding points produced by the associated MRG,
  in dimensions 2 to {\tt T}.  This is used only when $J > 1$.
 \endtab
\code

PROCEDURE EcrisGen (G : Generateur;    Cat : INTEGER);
\endcode
 \tab
  Writes the parameters and the results relative to a retained 
  generator {\tt G}, in category {\tt Cat}.
  Called by {\tt WriteSeekRes}.
 \endtab
\code

PROCEDURE WriteLatHead;
\endcode
 \tab
  Writes a heading for the results of a lattice test.
  Called by {\tt LatGen}.
 \endtab
\code

PROCEDURE WriteLatHeadTex;
\endcode
 \tab
  Writes a heading for the results of a lattice test in \LaTeX\ form.
  Called by {\tt LatGen}.
 \endtab
\code

PROCEDURE WriteDataComponent (j : Indexj);
\endcode
 \tab
  Write the statistical data relative to component {\tt j}.
 \endtab
\code

PROCEDURE WriteStatComponent (j : Indexj);
\endcode
 \tab
   Writes the statistics on the searches for primitive 
   (or examined) polynomials for the component {\tt j}.
   Called by {\tt WriteSeekRes}.
 \endtab
\code

PROCEDURE WriteTabBest (l : LONGINT;  G : Generateur);
\endcode
 \tab
  Prints a one-line summary of the figures of merit in each interval,
  for best generator retained in category (i.e., interval) number {\tt l}.
 \endtab
\code

PROCEDURE WriteSeekHead;
\endcode
 \tab
  Writes, in the current output file, the heading and the data 
  for a search by the {\tt SeekGen} procedure.
  Writes also the factorizations in the files if this was requested.
 \endtab
\code 

PROCEDURE WriteSeekRes;
\endcode
 \tab
   Writes the results of a search in the current output file,
   then closes it.
 \endtab
\code

PROCEDURE WriteSeekGen;
\endcode
 \tab
   Writes in a {\tt .gen} file the generators found in a search.
   If the data was read in {\tt xxx.dat}, then this procedure
   will open the file {\tt xxx.gen} and write in it.
 \endtab
\code

PROCEDURE SetupSeekGen;
\endcode
 \tab
  Called by {\tt SetUp}.
  Performs initializations for the searches of generators.
  Creates the ordered lists to store the `best' {\tt NbGen [C]}
  generators in each category {\tt C}.
  Initializes the bounds on the figures of merit in each dimension.
  Initializes the polynomial lists (to empty) and the search zones.
 \endtab
\code

PROCEDURE Setup;
\endcode
 \tab
  Obtains and opens data file, reads data and performs initializations.
  Initializes {\tt REDBAS} and {\tt MRG}.
  This procedure is called by {\tt SeekGen} and {\tt LatGen}.
 \endtab
\code

PROCEDURE LatGen;
\endcode
 \tab  Applies the required tests 
   (for ``lat'' programs) and prints results.
 \endtab
\code

PROCEDURE SeekGen;
\endcode
 \tab  Performs a seek for ``good'' generators 
   (implements ``seek'' programs).
 \endtab
\code

PROCEDURE CalcLaIndicesStreams (s, w : LONGINT);
\endcode
 \tab Calculates lacunary indices by groups of $s$ spaced by $2^w$
 \endtab 
\code

PROCEDURE LatGenStreams;
\endcode
 \tab  Applies the test for different values of $w$ where a big generator
  is divided into smaller ones, which are separated by $2^w$. 
 \endtab 
\code 

END LATIO.
\endcode









\defmodule{MRG}

This module offers tools to analyze simple or combined 
multiple recursive (linear congruential) generators (MRGs).
One must first initialize the module with a given modulus,
a given order, and a maximal dimension, by calling {\tt InitMRG}.
Each MRG will then be defined by a vector of multipliers {\tt a},
where {\tt a[i]} represents $a_i$ and the MRG is defined by
 $$ x_n := (a_1 x_{n-1} + \cdots + a_k x_{n-k}) \mod m.$$
One can then build the lattice basis associated to a vector of multipliers
for a given dimension, or apply the spectral test, or the Beyer test,
or both simultaneously, for all dimensions up to a specified limit.
The multipliers and modulus can be represented as {\tt LONGINT} or
{\tt SuperInteger}.  This can be changed by the command {\tt select}.

\bigskip\hrule
\code

DEFINITION MODULE MRG;

FROM MULT     IMPORT MScal, MVect;
FROM BASIS    IMPORT BVect;
FROM CONFIG   IMPORT MDim, BDim, OutputType;
FROM LATBASIS IMPORT Basis;
FROM TESTLAT  IMPORT TestConfig;
FROM NORM     IMPORT CriterionType, NormType, Normaliz;
FROM REDBAS   IMPORT VerifyType;
FROM SUPINT   IMPORT SuperInteger;
FROM SUPFACT  IMPORT Factors;
IMPORT BASISLR, MULTLI, MULTSI;

TYPE
   MRGen =  POINTER TO InfoMRG;
   InfoMRG = RECORD
      Bas   : Basis;
      k     : MDim;                     (* k = Order  of the recurrence *)
      aLI   : MULTLI.MVect;             (* Coefficients of the recurrence *)
      aSI   : MULTSI.MVect;             (*               "                *)
      END;



   LatticeType = (Full, Recurrent, Orbit, PrimePower);
\endcode
 \tabb Indicates which type of lattice is analyzed: that generated 
  by all possible states ({\tt Full}), 
  or by the set of recurrent states ({\tt Recurrent}), 
  or by a single forward orbit of the generator
  ({\tt Orbit}) for the case where it has several cycles, 
  or by a single orbit in the case where
  some component is an MLCG with modulus that is a power of a prime $p$
  ({\tt PrimePower}).  In the latter case, one must divide the modulus
  by an appropriate power of $p$, as described in \cite{rLEC97c},
  before calling a test procedure. 
  When {\tt LatticeType = Orbit}, the initial state must be specified
  in the array {\tt InSta} (see below).
  See also the data field {\em LatticeType\/} for the program {\tt seekl}.
 \endtabb
\code


VAR
   Bas : Basis;
\endcode
 \tabb  Lattice basis built and used by the procedure {\tt Test}.
   Can be recovered after calling {\tt Test}.
 \endtabb
\code

   Q2, D2, S2 : BASISLR.BVect;
\endcode
 \tabb  Those vectors contain the {\em squares\/} of the Beyer 
   quotients $q_t$, of the distances between
   hyperplanes $d_t$, and of the values of $S_t$, respectively, computed in
   the last call of {\tt Test} or {\tt TestLa}.
 \endtabb
\code

   Mins2, Minq2 : BASISLR.BVect;
\endcode
 \tabb  Gives the current minimal values of 
   $S_t$ and $q_t$ for a generator to be considered, in each dimension.
   The module initializes these bounds to zero.
 \endtabb
\code
 
   MaxAllDim : BOOLEAN;
\endcode
 \tabb  {\tt TRUE} if we are always looking for generators with the
  highest figures of merit, in all the dimensions.
  {\tt FALSE} if we want to {\em minimize\/} the figure of merit in
  some set of dimensions (that is, in some interval or category).
 \endtabb
\code

   TestCompleted : BOOLEAN;
\endcode
 \tabb  {\tt TRUE} if the last call to {\tt Test} or {\tt TestLa}
  has been successful, {\tt FALSE} otherwise.
 \endtabb
\code

   Lacunary : BOOLEAN;
\endcode
 \tabb  When this variable is {\tt TRUE}, the module treats lacunary
   indices; otherwise it assumes that the vectors are formed with
   successive indices.
 \endtabb
\code

   LatType : LatticeType;
\endcode
 \tabb  Indicates which lattice or sublattice is analyzed.
  See the data field {\em LatticeType\/} in {\tt seekl}.
  However, in the case {\tt PrimePower}, the modulus of each component
  must be divided appropriately {\em before\/} calling {\tt Test}
  or {\tt TestLa}; these test procedures assume that the modulus has 
  already been ajusted (reduced).
 \endtabb
\code

   DualLattice : BOOLEAN;
\endcode
 \tabb When {\tt TRUE}, we seek short vectors in the {\em dual\/}
  lattice (e.g., to compute the distance between adjacent hyperplanes).
  When {\tt FALSE}, we seek short vectors in the {\em primal\/}
  lattice.  The default value is {\tt TRUE}.
 \endtabb
\code

   UsedNorm : NormType;
\endcode
 \tabb  Which norm type is used to measure the length of the vectors.
  The default value is {\tt L2Norm}.
  WARNING: For other norms, the implementation is only partial, and
  only {\tt Spectral} is admitted for the type of test.
 \endtabb
\code

  NormalG : Normaliz;
\endcode
 \tabb To normalize the length of the shortest vectors (See module NORM).
 \endtabb
\code

   InvertCrit : BOOLEAN;
\endcode
 \tabb  If {\tt TRUE}, the value of the criterion $S_t$ is {\em inverted\/} 
  in the results, i.e., the criterion $S_t$ becomes a number larger than 1,
  the {\em smaller\/} the better.
 \endtabb
\code

   InSta : ARRAY MDim OF SuperInteger;    
\endcode
 \tabb If {\tt LatType = Orbit} then the user must specify a seed which
  will be stored in {\tt InSta}. 
 \endtabb
\code

PROCEDURE InitMRG (m : MScal;  k : MDim;  d : BDim);
\endcode
  \tab  From now on, this module will treat multiple recursive generators
   of order {\tt k} with modulus {\tt m}. 
   Dimensions of bases cannot exceed {\tt d}.
   (One can recall this procedure to change {\tt m}, {\tt k}, or {\tt d}).
   This procedure calls {\tt MULT.MSetm}.
\endtab
\code

PROCEDURE DeleteMRG (d : BDim);
\endcode
  \tab Deallocates the memory allocated to the matrices of {\tt SuperIntegers}
  {\tt VSI} and {\tt WSI}, and to the vectors of type {\tt BVect}, 
  {\tt xi, LacI}
  and {\tt q}, at a previous call to {\tt InitMRG}. It also calls 
  the procedure {\tt LATBASIS.DeleteBasis}, because {\tt LATBASIS.InitBasis} is
  called by {\tt InitMRG}. Must be called only if {\tt InitMRG} has been called
  previously.
\endtab
\code

PROCEDURE BuildBasis (B : Basis;   VAR a : MVect;   n : BDim);
\endcode
   \tab  Builds the basis directly in dimension {\tt n}.
   \endtab
\code

PROCEDURE IncDimBasis (B : Basis);
\endcode
   \tab  Increments by 1 the dimension of the basis.
   \endtab
\code

PROCEDURE Test
   ( VAR a          : MVect;
     T              : CriterionType;
     FromDim, ToDim : BDim;
     VerifyBB       : VerifyType;
     Out            : OutputType );

\endcode
 \tab  Applies the test {\tt T} (Beyer and/or Spectral)
   to the MRG generator corresponding to {\tt a}, in dimensions
   {\tt FromDim} to {\tt ToDim}.
   The initial basis is built automatically, and during the procedure,
   the variable {\tt Bas} will contain the current basis.
   {\tt VerifyBB} indicates if
   the results must be verified in a second pass, with error bounds
   on the floating-point computations for the branch-and-bound.
   The parameter {\tt Out} specifies the form of output produced.
   If a test fails, then the procedure stops and the variable
   {\tt TestCompleted} is set to {\tt FALSE}.
 \endtab
\code

PROCEDURE TestCri
   ( VAR a          : MVect;
     T              : CriterionType;
     FromDim, ToDim : BDim;
     VerifyBB       : VerifyType;
     Out            : OutputType;
     r              : BDim;
     n              : INTEGER );
\endcode
  \tab Does the same thing as {\tt Test}, but two additional parameters
  {\tt r} and {\tt n} must be specified, which are the rank and the copy 
  factor for copy-rules, and then the basis are modified accordingly.
  Must be used only if {\tt FromDim} is larger than {\tt r}. Not very useful
  in the context of LCG's, but it would have been complicated to write
  this procedure and the one it calls, {\tt CriSpectralTest()}, in a
  separate module, because they use variables internally declared in 
  {\tt MRG}. The LCG corresponding to the rank 1 rule must have its
  modulo and multiplier modified before {\tt TestCri} is called.
  THIS PROCEDURE IS NOT IMPLEMENTED. 
  \endtab
\code

PROCEDURE BuildLaBasis
   (B : Basis;   VAR a : MVect;   VAR I : BVect;   n, MaxDim : BDim);
\endcode
 \tab  Builds the (lacunary) basis directly in dimension {\tt n}.
   One must indicate MaxDim and use IncDimLaBasis only for dimensions
   smaller than this.
   The vector {\tt I} must contain the {\tt MaxDim} lacunary indices.
 \endtab
\code

PROCEDURE IncDimLaBasis (B : Basis);
\endcode
  \tab  Increments by 1 the dimension of the basis. The basis
   must have been initialized by {\tt InitLaBasis}
   and the dimension of {\tt B} must be less than {\tt MaxDim}.
  \endtab
\code

PROCEDURE TestLa
   ( VAR a          : MVect;
     VAR I          : BVect;
     T              : CriterionType;
     FromDim, ToDim : BDim;
     VerifyBB       : VerifyType;
     Out            : OutputType );
\endcode
 \tab  Similar as {\tt Test}, but with lacunary indices.
  The vector {\tt I} must contain the {\tt ToDim} lacunary indices.
 \endtab
\code

PROCEDURE PerMax (VAR a : MVect;   VAR ListmMS1, Listr : Factors) : BOOLEAN;
\endcode
  \tab  Checks if the generator corresponding to {\tt a}
    has maximal period. If {\tt m} is not prime then the procedures stops
    and {\tt TestCompleted} is set to {\tt FALSE}.
    {\tt ListmMS1} and {\tt Listr} must contain the factors of
    {\tt m-1} and {\tt r}, or {\tt NIL} if the factors are unknown.
    (See {\tt MULT.SetFactors}.)
   \endtab
\code

END MRG.
\endcode


\defmodule{TESTLAT}

This module offers a procedure to test several projections of the lattice specified by a modulus 
and  a set of generating vectors.  Several categories of projections can be considered, with 
specific projections in each category.  All these parameters are specified via 
{\tt InitTestConfig}, stored in a {\tt TestConfig} structure which is passed to the
{\tt TestLattice} procedure.


\bigskip\hrule
\code

DEFINITION MODULE TESTLAT;

FROM BASIS    IMPORT BVect;
FROM MULT     IMPORT MScal, MVect;
FROM CONFIG   IMPORT MDim, BDim, DimProj, MCoord, Categ, OutputType;
FROM NORM     IMPORT CriterionType, Normaliz, NormType, GammaType;
FROM LATBASIS IMPORT Basis;
FROM REDBAS   IMPORT VerifyType;
IMPORT BASISLR;


TYPE
   CoordSet = SET OF MCoord;

   TestConfig = POINTER TO InfoTestConfig;
   InfoTestConfig = RECORD
      m             : MScal;
\endcode
\tabbb
  The modulus.
\endtabbb
\code

      k             : MDim;
\endcode
\tabbb
  Number of generating vectors (usually equal to the order of the
  linear recurrence).
\endtabbb
\code

      NbCoord       : MCoord;
\endcode
\tabbb
  Total number of coordinates that are considered.
\endtabbb
\code

      Maxds         : BDim;
      Maxdp         : DimProj;
\endcode
\tabbb
  {\tt Maxds} is the maximum dimension of the projections over successive 
  coordinates.
  {\tt Maxdp} is the maximum dimension of the projections over nonsuccessive
  coordinates.
  These are the values used for memory allocation in {\tt InitTestConfig}.
\endtabbb
\code

      NbCats, C     : Categ;
\endcode
\tabbb
  Number of categories (or lists) that are considered.
  The variable {\tt C} is initially set to {\tt NbCats} by {\tt TestLattice},
  and is decreased when the generator is not good enough for the 
  category currently considered.
\endtabbb
\code

      DimenProj     : ARRAY Categ OF BDim;
      RangeProj     : ARRAY Categ OF ARRAY DimProj OF BDim;
      DimRegLag     : ARRAY Categ OF BDim;
      RangeRegLag   : ARRAY Categ OF ARRAY DimProj OF BDim;
\endcode
\tabbb
   These arrays specify which projections are considered for each category.
   {\tt DimenProj[c]} is the maximum number of dimensions for which the
   projections over non-successive and non-regular indices are considered.
   (E.g., if the value is 3, the pairs and triples are considered.)
   For $j>1$, {\tt RangeProj[c,j]} indicates the maximal value of $i_j+1$
   that is considered for $j$-dimensional projections in category $c$,
   assuming that $i_1=0$.  For $j=1$, it gives the maximum dimension for
   the successive dimensions, for category $c$.
   Non-successive indices that are regularly spaced, i.e., where
   $i_j = (j-1)s$ for some $s$, are also considered.
   For category $c$, {\tt DimRegLag[c]} gives the maximum value of $s$
   and {\tt RangeRegLag[c,s]} gives the maximum value of $j$ (number
   of dimensions) that is considered for this $s$.
\endtabbb
\code

      Mins2, Maxs2  : ARRAY Categ OF LONGREAL;
\endcode
\tabbb
   For each category $c$, gives the minimum and maximum value of the 
   worst-case criterion within category $c$ for a generator to be 
   considered for category $c$ and the categories above it.
   Whenever {\tt ValCrit[c]} is outside the range
   [{\tt Mins2[c], Maxs2[c]}], the corresponding generator need not be
   considered any longer for all the categories $c'\ge c$.
   It may still be considered, however, for the categories $c'<c$.
\endtabbb
\code

      InvertL, InvertS : BOOLEAN;
\endcode
\tabbb
   Indicates if we want to print the multiplicative inverses of
   the length of the shortest vector and/or the figure of merit {\tt S2}.
   For the spectral test with the Euclidean norm, for instance, 
   if {\tt InvertL = TRUE}, we obtain the distance between hyperplanes
   instead of the length of the shortest vector.
   Default values are {\tt FALSE}.
\endtabbb
\code

      Criterion     : CriterionType;
\endcode
\tabbb
   Indicates which criterion is used. 
\endtabbb
\code

      Normal        : Normaliz;
\endcode
\tabbb
   {\tt Normal} specifies the normalization method
   used to compute {\tt S2}.
\endtabbb
\code

      Verify        : VerifyType;
\endcode
\tabbb
   Indicates the level of verification (i.e., {\tt REDBAS2} or
   {\tt REDBAS3} is used).
\endtabbb
\code

      Output        : OutputType;
\endcode
\tabbb
   Indicates what kind of output is written during the test.
\endtabbb
\code

      MasterGen     : ARRAY MDim OF ARRAY MCoord OF MScal;
\endcode
\tabbb
  {\tt MasterGen} contains an initial set of $k$ generating vectors
  for all the {\tt NbCoord} coordinates that are to be considered in the 
  projections.
  The projections of these vectors, obtained by selecting the 
  appropriate columns in {\tt MasterGen}, together with additional vectors
  of the form $(\dots, m, \dots)$, are used to build the triangular
  bases in {\tt TBas}.
\endtabbb
\code

      TBas          : Basis;
\endcode
\tabbb
  Contains a triangular version of the basis
  for the {\tt d} coordinates that are currently under consideration.
  Used by the {\tt TestProjections} procedure.
\endtabbb
\code

      RBas          : ARRAY DimProj OF Basis;
\endcode
\tabbb
  {\tt RBas[d]} contains a reduced basis
  for the {\tt d} coordinates that are currently under consideration.
  This is used by the {\tt TestProjections} procedure, which goes
  back and forth in this array while exploring the tree of projections.
  {\tt RBas[0]} is used for the lattices that correspond to successive
  coordinates.
  
\endtabbb
\code

      L2, S2        : ARRAY DimProj OF BASISLR.BVect;
\endcode
\tabbb
   The vector {\tt L2} will contain the lengths of the shortest vectors 
   raised to the power $p$ in case of the spectral test with the 
   $\cL_p$ norm.  More precisely, {\tt L2[1,t]} is the value for the
   lattice that corresponds to successive coordinates in $t$ dimensions,
   whereas {\tt L2[d,t]}, for $d>1$, is the worst value for all the 
   $d$-dimensional projections over coordinate sets $I = \{i_1,\dots,i_d\}$
   where $i_1=0$ and $i_d=t-1$.  In all cases,   
   {\tt S2[d,t]} will contain the corresponding normalized quantity,
   also raised to the power $p$, or the Beyer quotient in case of the 
   Beyer test.
\endtabbb
\code

      ValCrit : ARRAY Categ OF LONGREAL;
\endcode
\tabbb
   {\tt ValCrit[c]} will contain the worst-case value of the figure 
   of merit in {\tt S2[d,t]} over the set of projections that are considered
   for category {\tt c}.
\endtabbb
\code

      OK : BOOLEAN;
\endcode
\tabb
   Indicates if the last operation (e.g., call to {\tt TestLattice,
   SpectralTest}, etc.) has been successful.
\endtabb
\code
      END;



PROCEDURE CreateTestConfig (VAR T : TestConfig);
\endcode
  \tab  
  Creates the structure {\tt T}.  One must then call
  {\tt InitTestConfig} to allocate memory for the elements of
  {\tt MasterGen, TBas, RBas, L2}, and {\tt S2}.
  \endtab
\code

PROCEDURE DeleteTestConfig (VAR T : TestConfig);
\endcode
  \tab  
  Deletes the structure {\tt T}, after deallocating the memory
  allocated to its fields by {\tt InitTestConfig}.
  \endtab
\code

PROCEDURE InitTestConfig 
  ( T : TestConfig;  Norm : NormType;
    m  : MScal;   k  : MDim;      NbCoord : MCoord;  
    ds : BDim;    dp : DimProj;   NbCats  : Categ;
    gamma : GammaType;   beta : LONGREAL
  );
\endcode
  \tab
  Initializes the structure {\tt T}.
  In the parameters, {\tt m} represent the modulus, {\tt k} the number of
  generating vectors, 
  {\tt NbCoord} the total number of coordinates to be considered,
  {\tt ds} the maximum dimension of the projections over successive 
  coordinates,
  {\tt dp} the maximum dimension of the projections over nonsuccessive
  coordinates,
  and {\tt NbCats} the number of categories.
  This procedure creates the elements of {\tt MasterGen[1..k, 1..NbCoord]},
  {\tt L2}, and {\tt S2},
  creates the bases {\tt TBas} and {\tt RBas[0..d]},
  initializes their sizes appropriately, their modulus to $m$,
  and their {\tt Norm} field to {\tt Norm}.
  The parameters {\tt gamma {\rm and} beta} are used to initialize {\tt
  T\^{}.Normal} by calling {\tt InitNormaliz}.
  \endtab
\code

PROCEDURE AddToBasis (T : TestConfig;  d : BDim;  Ind : CoordSet;  n : BDim);
\endcode
  \tab
  Assumes that {\tt MasterGen} has been initialized to a proper set
  of generating vectors.
  Adds a vector to the basis {\tt TBas} that corresponds to adding 
  the $j$th coordinate of the generating vectors, and updates the dual.
  \endtab
\code

PROCEDURE RemoveFromBasis (T : TestConfig; d : BDim;  Ind : CoordSet; n : BDim);
\endcode
  \tab
  Removes from the basis {\tt TBas} the vector that corresponds to the 
  $j$th coordinate, and updates the dual.
  \endtab
\code

PROCEDURE TestLattice (T : TestConfig);
\endcode
 \tab  Applies the tests specified by {\tt T} to the lattice
   whose set of generating vectors is in its field \verb=MasterGen=.
 \endtab
\code

END TESTLAT.
\endcode

\defmodule {REDBAS}

This module offers tools to reduce a basis in various ways,
and to find a shortest vector in a lattice
using pre-reductions and a branch-and-bound (BB) procedure.
It also has a procedure to compute a Minkowski-reduced basis
(this is much more expensive than finding a shortest vector).
%%  See \cite{rCOU96x} for more details.
One must first call {\tt InitREDBAS} before calling any other procedure
in this module.

\bigskip\hrule
\code

DEFINITION MODULE REDBAS;

IMPORT BASISLR;
FROM LATBASIS IMPORT Basis;
FROM CONFIG   IMPORT BDim;

TYPE
   VerifyType  = (Verify0, Verify2, Verify3);
\endcode
 \tabb Type of verification performed by the tests:
    uses {\tt REDBAS},  {\tt REDBAS2}, or {\tt REDBAS3}
    for {\tt Verify0}, {\tt Verify2}, and {\tt Verify3},
   respectively.
 \endtabb
\code


VAR
   PreRedDieterSV, PreRedLLLSV, PreRedLLLRM : BOOLEAN;
\endcode
  \tabb  These boolean variables indicate which type of pre-reduction 
   is to be performed for {\tt ShortestVector} (SV) and for 
   {\tt ReductMinkowski} (RM).
   {\tt Dieter} means the pairwise pre-reduction as in the procedure
   {\tt PreRedDieter}.  
   {\tt LLL} means the LLL reduction of Lenstra, Lenstra, and Lov\'asz.

   The variable {\tt PreRedDieterSV} is originally set to {\tt TRUE}
   and the 2 others are originally set to {\tt FALSE}.
   These variables are reset automatically depending on the thresholds 
   {\tt MinkLLL, ShortDiet, ShortLLL} as explained below.
   Note that the LLL reduction is never performed when {\tt SISquares = TRUE}.
   \pierre {Reason?}
 \endtabb
\code

   ShortDiet : LONGINT;
\endcode
 \tabb
   Whenever the number of nodes in the BB tree exceeds the threshold 
   {\tt ShortDiet}, in the {\tt ShortestVector} procedure, {\tt PreRedDieterSV}
   is automatically set to {\tt TRUE} for the next call; 
   otherwise it is set to {\tt FALSE}.
   The default value is 1000.
 \endtabb
\code

   ShortLLL : LONGINT;
\endcode
 \tabb 
   Whenever the number of nodes in the BB tree exceeds the threshold 
   {\tt ShortLLL}, in the {\tt ShortestVector} 
   procedure, {\tt PreRedLLLSV}
   is automatically set to {\tt TRUE} for the next call; 
   otherwise it is set to {\tt FALSE}.
   The default value is 1000.
  \endtabb
\code

   MinkLLL : LONGINT;
\endcode
 \tabb  
   Whenever the number of nodes in the BB tree exceeds {\tt MinkLLL}
   in the {\tt ReductMinkowski} procedure,
   {\tt PreRedLLLRM} is automatically set to {\tt TRUE} for the
   next call; otherwise it is set to {\tt FALSE}.
   The default value is 500\,000.
  \endtabb
\code

   MaxPreRed : LONGINT;
\endcode
 \tabb
   Maximum number of transformations in the procedure {\tt PreRedDieter}.
   After {\tt MaxPreRed} successful transformations have been performed,
   the prereduction is stopped.
   The default value is 1\,000\,000.
  \endtabb
\code

   MaxNodesBB : LONGINT;     
\endcode
  \tabb The maximum number of nodes that we are ready to accept
   in the branch-and-bound (BB) tree when calling {\tt ShortestVector}
   or {\tt ReductMinkowski}.  When this number is exceeded, the procedure
   aborts and returns {\tt FALSE}.
   The default value is 10\,000\,000.
  \endtabb
\code


PROCEDURE InitREDBAS (MaxDim : BDim);
\endcode
  \tab  Initializes the internal variables of this 
   module with a given upper dimension {\tt MaxDim}
   for the bases, for memory allocation purposes.
   This {\tt MaxDim} (saved in the internal variable {\tt CurMaxDim})
   should be at least as large as the dimension of any
   {\tt Basis B} that this module has to work on (i.e., {\tt B\^{}.Dim}
   can never exceed {\tt CurMaxDim}).
   \endtab
\code

PROCEDURE DeleteREDBAS;
\endcode
  \tab Deallocates the memory which was allocated when calling
  {\tt InitREDBAS} for the last time.
% for the vector of {\tt SuperIntegers} 
%  {\tt Sv}, the matrices {\tt VTemp}, and 
%  {\tt WTemp}, and the vector {\tt ZeroVect}.
   Must be called only if {\tt InitREDBAS} has been called previously.
  \endtab
\code

PROCEDURE GetMaxDim (): BDim;
\endcode
  \tab  Returns the current value of {\tt CurMaxDim}.
   \endtab
\code

PROCEDURE PairwiseRedPrimal (B : Basis;  i, d : BDim);
\endcode
  \tab  Performs pairwise reductions.
  This procedure tries to reduce each basis vector with 
  index larger than $d$, and distinct from $i$, by adding to it a 
  multiple of the $i$-th vector.  
  Always use the Euclidean norm.
 \endtab
\code

PROCEDURE PairwiseRedDual (B : Basis;  i : BDim);
\endcode
  \tab  Performs pairwise reductions, trying to reduce every other vector of
  the {\em dual\/} basis by adding multiples of the $i$-th vector.
  That may change the $i$-th vector in the primal basis.
  Each such dual reduction is actually performed only if that does not
  increase the length of vector $i$ in the primal basis. 
  Always use the Euclidean norm.
 \endtab
\code

PROCEDURE PreRedDieter (B : Basis;  d : BDim);
\endcode
  \tab  Performs the reductions of the preceding two procedures 
  using cyclically all values of $i$ (only for $i > d$ in the latter case)
  and stops after either {\tt MaxPreRed} successful transformations have been 
  achieved or no further reduction is possible.
  Always use the Euclidean norm.
  \endtab
\code

PROCEDURE SetScale (B : Basis);
\endcode
\tab  This procedure is used when {\tt SISquares = TRUE}, 
   just before doing the Choleski decomposition
   for the branch-and-bound procedure, in the procedures {\tt ShortestVector}
   and {\tt ReductMinkowski}.
   It rescales the square Euclidean lengths of the basis and $m$-dual 
   basis vectors by factors {\tt EchVV} and {\tt EchWW}, respectively,
   and $m^2$ by the product {\tt EchVV * EchWW}.
   These factors are chosen by the procedure 
   so that the log of the largest squared
   Euclidean length is not too far from 0.
   The rescaled values are placed in {\tt VVLR} and {\tt WWLR}, and the 
   rescaled value of $m^2$ in {\tt mLR2}.
   The Choleski decomposition and the branch-and-bound procedures
   use these rescaled values to avoid {\tt LONGREAL} overflow.
\endtab
\code

PROCEDURE RedLLL (B : Basis;  fact : LONGREAL;  m : LONGINT;  Max : BDim);
\endcode
  \tab  Performs a LLL basis reduction with coefficient {\tt fact}, which must
   be smaller than one.  If {\tt fact} is closer to one, the basis will be
   (typically) ``more reduced'', but that will require (slightly) more work.
   Always use the Euclidean norm.
   {\em Danger: With the current implementation, there could be {\tt LONGREAL}
   overflow when {\tt SISquares = TRUE} and there are very long basis vectors.}
  \endtab
\code

PROCEDURE CalculCholeski (B : Basis;  VAR DC2 : BASISLR.BVect;
                          VAR C0 : BASISLR.BMatr) : BOOLEAN;
\endcode
\tab  Performs a Choleski decomposition of the matrix of scalar
  products of the basis vectors, i.e., $\bU'\bU = \bV'\bV$ where $\bV$
  is the matrix of the basis vectors and $\bU$ is upper triangular.
  Returns in {\tt C0} the elements of $\bU$ that are {\em strictly above\/}
  the diagonal, i.e., $u_{i,j}$ for $j > i$.
  Returns in {\tt DC2} the squared elements of the diagonal of $\bU$.
  If {\tt SISquares = TRUE}, these elements are rescaled (divided) by
  {\tt EchVV}.  Returns {\tt TRUE} if the decomposition was successful.
\endtab
\code

PROCEDURE TransformStage3 (B : Basis;  VAR z : BASISLR.VectLI;  VAR k : BDim);
\endcode
 \tab Procedure used in {\tt ReductMinkowski} to perform a transformation
  of stage 3 described in \cite{rAFF85a}. Also used in {\tt ShortestVector}.
  Assumes that $\sum_{i=1}^t z_i V_i$ is a short vector that will enter
  the basis.  Tries to reduce some vectors by looking for indices $i < j$
  such that $|z_j| > 1$ and $q=\lfloor z_i/z_j\rfloor \not=0$,
  and adding $q V_i$ to $V_j$ when this happens.
  Returns in $k$ the last index $j$ such that $|z_j|=1$.
 \endtab
\code

PROCEDURE ShortestVector (B : Basis) : BOOLEAN;
\endcode
  \tab Computes the shortest vector in lattice with basis {\tt B},
   using branch-and-bound.
   If {\tt MaxNodesBB} is exceeded during {\em one} of the branch-and-bounds,
   the procedure aborts and returns {\tt FALSE}.
   Otherwise, it returns {\tt TRUE}.
%  This procedure does not care about numerical imprecision, and therefore,
%  the results are not 100\% reliable.
\endtab
\code

PROCEDURE ReductMinkowski (B : Basis;  d : BDim) : BOOLEAN;
\endcode
   \tab  Reduces {\tt B} into a Minkowski reduced basis, with respect to the
   Euclidean norm, assuming
   that the first $d$ vectors are already reduced and sorted.
   If {\tt MaxNodesBB} is exceeded during {\em one} of the branch-and-bounds,
   the procedure aborts and returns {\tt FALSE}.
   Otherwise, it returns {\tt TRUE}, and the basis is reduced and
   is sorted by vector lengths (the shortest vector is {\tt V[1]} and
   the longest is {\tt V[Dim]}).
   This procedure does not care about numerical imprecision due to the
   (64-bit) floating-point representation.
   In this sense, the results are not 100\% reliable.
   \endtab
\code

END REDBAS.
\endcode



\defmodule {REDBAS2}

This module offers tools to reduce a basis in the
sense of Minkowski and find a shortest vector in a lattice.
The difference between this and {\tt REDBAS} is that here,
bounds on the numerical errors are computed to make sure that the
results are formally correct.

\bigskip\hrule
\code

DEFINITION MODULE REDBAS2;

FROM CONFIG   IMPORT BDim;
FROM LATBASIS IMPORT Basis;

PROCEDURE InitREDBAS2 (d : BDim);
\endcode
   \tab  Initializes the module. The value of {\tt d} should be at least
    as large as the dimension of the largest basis to be considered
    by the procedures below. \endtab
\code

PROCEDURE VerifMinkowski (B : Basis; d : BDim) : BOOLEAN;
\endcode
   \tab  Performs a formal verification and 
   returns {\tt TRUE} if the basis {\tt B} is really reduced. 
   Assumes that the first {\tt d} vectors are really reduced and
   verifies only the remaining ones. \endtab
\code

PROCEDURE ReductVerifMinkowski (B : Basis) : BOOLEAN;
\endcode
   \tab  Performs a reduction as in {\tt ReductMinkowski}, then verifies it
   formally, and if the basis is not really reduced, completes the reduction.
   This is usually much more costly than {\tt ReductMinkowski}, but the
   result is perfectly reliable. \endtab
\code

PROCEDURE VerifShortest (B : Basis) : BOOLEAN;
\endcode
   \tab  Performs a formal verification and 
   returns {\tt TRUE} if the first vector of basis {\tt B} 
   is the shortest vector in the lattice generated by {\tt B}.
   \endtab
\code

PROCEDURE ShortestVectorVerif (B : Basis) : BOOLEAN;
\endcode
   \tab  Computes the shortest vector as in {\tt ShortestVector},
   performs a formal verification, and if the vector was not really the 
   shortest due to numerical (floating-point) imprecisions, 
   computes the shortest one.
   Returns {\tt TRUE} if the procedure succeeds.
   This is usually much more costly than {\tt ShortestVector}, but the
   result is perfectly reliable.  \endtab
\code

END REDBAS2.
\endcode

\defmodule{REDBAS3}

This module offers tools similar to those of {\tt REDBAS2},
except that the ``provably correct'' results are obtained by a
different method than in {\tt REDBAS2}, bypassing all floating-point
calculations, as explained in \cite{rCOU96x}.

\bigskip\hrule
\code

DEFINITION MODULE REDBAS3;

FROM CONFIG   IMPORT BDim;
FROM LATBASIS IMPORT Basis;
FROM BASISLR  IMPORT VectLI;

PROCEDURE VerifMinkowski (B : Basis;  d : BDim) : BOOLEAN;
\endcode
  \tab  Performs a formal verification and 
   returns {\tt TRUE} if the basis {\tt B} is really reduced. \endtab
\code
\hide (*
PROCEDURE ReductVerifMinkowski (B : Basis) : BOOLEAN;
*)  \endhide

PROCEDURE VerifShortest (B : Basis; d1, d2 : BDim; VAR z : VectLI) : BOOLEAN;
\endcode
   \tab  Performs a formal verification and
   returns {\tt TRUE} if the first vector of basis {\tt B}
   is the shortest vector in the lattice generated by the first {\tt d1}
   vectors of {\tt B}.   The parameter {\tt d2} is a threshold,
   which should be between 1 and the dimension of the lattice,
   and may affect the performance.  To compute the relevant
   coefficients (expressed as determinants in the primal or dual basis),
   we use the primal basis up to dimension {\tt d2}, and then the
   dual basis.   When the procedure returns {\tt FALSE}, then the
   coefficients of a shorter vector with respect to the basis are
   returned in {\tt z}.
   \endtab
\code\hide
(*
PROCEDURE ShortestVectorVerif (B : Basis) : BOOLEAN;
*)
\endhide

END REDBAS3.
\endcode

\defmodule {LATBASIS}

This module offers tools to manipulate lattice bases
(see Section 1.3).
Each lattice is represented by a basis $V$ and its dual $W$.
It is sometimes possible, as in the case with lattices associated with
LCGs or MRGs, to multiply a lattice (and its dual) by a constant factor
in such a way that they are included in $\Z^t$, allowing exact representation
of basis vector coordinates.
The duality relation will now read $V_i\cdot W_j = m\delta_{ij}$
for some integer constant {$m$}.

\iffalse  %%%%%%%%%%%%%
There is a choice between two representations for the vector coordinates:
{\tt LONGREAL} or {\tt SuperInteger}.
The latter is a type defined in a package named SENTIERS which
offers tools for manipulating arbitrary large integers.
The {\tt SuperInteger} representation should be used
when the coordinates of the (dual) basis vectors, are too large
to fit a {\tt LONGREAL} without loss of precision.
But computations are then much slower.
The program {\tt select} decides which one is used.
All computations will be performed using appropriate procedures,
depending on the representation.
\fi  %%%%%%%%%%%%%%%%%

The squared lengths of the basis vectors and that of their duals
are stored as {\tt LONGREAL} in {\tt VVLR} and {\tt WWLR} respectively.
However, if they are too large to fit in a {\tt LONGREAL},
even approximately, then {\tt SISquares} is set to {\tt TRUE} and
they are stored as {\tt SuperInteger} in {\tt VVSI} and {\tt WWSI}.  

A variable of type {\tt Basis} is created by calling {\tt CreateBasis}.
The squared vector lengths {\tt VVLR} and {\tt WWLR} 
(or {\tt VVSI} and {\tt WWSI})
can then be updated by calling {\tt UpdateVVWW}.
But it is not necessary to update them before calling any of the procedures
of the modules {\tt REDBAS} and {\tt REDBAS2}
(including {\tt ReductMinkowski} and {\tt ShortestVector}).

\bigskip\hrule
\code

DEFINITION MODULE LATBASIS;

FROM SUPINT  IMPORT SuperInteger;

IMPORT BASISLR, BASISSI;
FROM BASIS    IMPORT BMatr, BScal;
FROM CONFIG   IMPORT BDim;
FROM NORM     IMPORT NormType;
FROM MULT     IMPORT MScal;

TYPE
   Basis = POINTER TO InfoBasis;
   InfoBasis = RECORD
      ResDim     : BDim;             (* Dimension of allocated memory.      *)
      Dim        : BDim;             (* Actual dimension.                   *)
      mLR        : LONGREAL;         (* Modulus m.                          *)
      mSI        : SuperInteger;     (* Modulus m.                          *)
      mLR2       : LONGREAL;         (* m^2, rescaled if SISquares.         *)
      mSI2       : SuperInteger;     (* Value of m^2 when SISquares.        *)
      V, W       : BMatr;            (* Basis and dual basis vectors.       *)
      Norm       : NormType;
      Lmin, Lmin2  : LONGREAL;       (* Length and squared length of shor-  *)
                                     (* test basis vector, using Norm.      *)
                                     (* Different meaning in REDBB.         *)
      SISquares    : BOOLEAN;
      VVLR, WWLR   : BASISLR.BVect;  (* Squared Euclidean vector lengths,   *)
                                     (* if not SISquares.                   *)
      EchVV, EchWW : SuperInteger;   (* Scaling factors when SISquares.     *)
      VVSI, WWSI : BASISSI.BVect;    (* Rescaled squared Euclidean vector   *)
                                     (* lengths, if SISquares.              *)
      XX : ARRAY BDim OF BOOLEAN;    (* Used for ReductMinkowski only.      *)
                                     (* Ask R. Couture for details.         *)
      END;
\endcode
  \tabb
   A lattice basis.  
   {\tt ResDim} is the dimension of the arrays allocated for the basis
   (the maximum value that {\tt Dim} can take).
   The memory for the {\tt BVect} and {\tt BMatr} elements is allocated
   for dimensions 0 to {\tt ResDim} if {\tt ResDim} $> 0$, whereas
   no memory is allocated if {\tt ResDim} $= 0$.

   {\tt Dim} is the current dimension of the basis,
   {\tt V} contains the basis, and {\tt W} the $m$-dual basis. 
   The integer constant $m$ used in the duality relation is stored
   as a {\tt LONGREAL} in {\tt mLR}, and as a {\tt SuperInteger}
   in {\tt mSI}.

   {\tt Norm} is the type of norm used to measure the vector lengths.
   {\tt Lmin} and {\tt Lmin2} are used to store the length and the 
   squared length of the shortest vector found to date, using this norm.
   (Usually, only one of these 2 variables is used at a time; 
   see {\tt REDBAS}.)

   The boolean {\tt SISquares} is set to {\tt TRUE} by {\tt InitBasis}
   when $m$ exceeds approximately $2^{480}$, and to {\tt FALSE} otherwise.

   When {\tt SISquares = FALSE}, {\tt mLR2} contains the value of $m^2$,
   {\tt VVLR} and {\tt WWLR} {\em may\/} memorize the squared Euclidean 
   lengths of the vectors in $V$ and in $W$. 
   They are not always updated, and when they are not, they
   are set to negative values.

   When {\tt SISquares = TRUE}, {\tt mSI2} contains $m^2$,
   {\tt mLR2} contains $m^2$ divided (rescaled) by {\tt EchVV * EchWW},
   the squared Euclidean vector lengths (again, not always updated),
   are stored in {\tt VVSI} and {\tt WWSI}, and 
   {\tt Lmin2} is rescaled by {\tt EchVV}.
   Also, the squared vector lengths rescaled by the factors
   {\tt EchVV} and {\tt EchWW}, are put in {\tt VVLR} and {\tt WWLR},
   respectively, by {\tt SetScale}.
   The scaling factors {\tt EchVV} and {\tt EchWW} are computed
   by the procedure {\tt SetScale} and set to values such that
   the log of the squared Euclidean length of the longest vector is 
   not too far from 0.  
   The aim of this rescaling is to avoid {\tt LONGREAL} overflow.
   \endtabb
%\hide
%VAR
%   (* This has been moved to  BASIS^.XX  in LATBASIS.  *)
%   XX  : ARRAY BDim OF BOOLEAN;    (* ??? *)    
%\endhide
\code

PROCEDURE CreateBasis (VAR B : Basis);
\endcode
 \tab Creates the data structure {\tt B} for a basis. 
    {\tt SISquares} is initialized to {\tt FALSE},
    {\tt Norm\/} is initialized to {\tt L2Norm},
    and {\tt ResDim} is set to 0.
    One must then call {\tt InitBasis} to reserve the memory for
    the {\tt BVect} and {\tt BMatr} elements.
 \endtab
\code

PROCEDURE InitBasis (B : Basis;  m : MScal;  Nor : NormType;  d : BDim);
\endcode
   \tab
    Initializes the modulus and the norm for the basis {\tt B}.  
    Reserves the memory for the {\tt BVect} and {\tt BMatr}
    fields of {\tt B}, for indices 0 to $d$ if $d>0$.
    Here, {\tt d} is the new value of {\tt ResDim}.
    If {\tt d < B\^{}.ResDim}, then the memory already reserved which
    is no longer needed is deallocated.
    \endtab
\code

PROCEDURE DeleteBasis (B : Basis);
\endcode
   \tab  Does the same as {\tt InitBasis (B, B\^{}.m, B\^{}.Norm, 0)}, 
    and then the  opposite of {\tt CreateBasis (B)}.
%   Deallocates the memory allocated for {\tt V}, {\tt W}, {\tt VVSI} 
%   and {\tt WWSI}, at a previous call to {\tt InitBasis}. 
%   Must be called only if {\tt InitBasis} has been called previously.
   \endtab
\code

PROCEDURE WriteBasis (B : Basis);
\endcode
  \tab  Writes content of {\tt B} in current output file.
  \endtab
\code

PROCEDURE ReadBasis (VAR B : Basis);
\endcode
  \tab  Reads a basis from current input file.
        The basis must have been written to the file by {\tt WriteBasis}.
  \endtab
\code

PROCEDURE CopyBasis (B1, B2 : Basis);
\endcode
 \tab  Copies {\tt B1} into {\tt B2}.  \endtab
\code

PROCEDURE Permute (B : Basis;  i, j : BDim);
\endcode
   \tab Exchanges vectors i and j in the basis and its dual. \endtab
\code

PROCEDURE SortBasis (B : Basis;  d : BDim);
\endcode
   \tab  Sorts the basis vectors (with index from $d+1$ to {\tt B\^{}.Dim}) 
   by increasing length.
   The dual vectors are permuted accordingly.
   Assumes that {\tt VVLR} and {\tt WWLR} 
   (or {\tt VVSI} and {\tt WWSI}, if {\tt SISquares}) are up to date.
%  Note: no longer updates {\tt VVLR} (or {\tt VVSI}). 
 \endtab
\code

PROCEDURE Dualize (B : Basis);
\endcode
  \tab Exchanges {\tt V, VVLR, VVSI, EchVV} and {\tt W, WWLR, WWSI,
       EchWW}.
 \endtab
\code

PROCEDURE UpdateVVWW (B : Basis;  d : BDim);
\endcode
  \tab  Updates {\tt VVLR} and {\tt WWLR} (or {\tt VVSI} and {\tt WWSI}
     if {\tt SISquares}) as the square Euclidean lengths of the vectors
     in {\tt V} and {\tt W} (for vectors 
     with index from $d+1$ to {\tt B\^{}.Dim}).
 \endtab
\code

PROCEDURE EraseVVWW (B : Basis);
\endcode
  \tab
    Sets the vectors {\tt VVLR} and {\tt WWLR} (or {\tt VVSI} and {\tt WWSI}
    if {\tt SISquares}) to negative values. 
\endtab
\code

PROCEDURE CheckDuality (B : Basis): BOOLEAN;
\endcode
  \tab
   Check that {\tt Basis B} satisfies the duality relation
   $V[i]\cdot W[j] = m\,\delta_{ij}$. If so, returns {\tt TRUE}, otherwise 
   returns {\tt FALSE}.
\endtab
\code

PROCEDURE BaseEquivalence (B1, B2 : Basis) : BOOLEAN;
\endcode
  \tab
   Check that {\tt Basis B1} and {\tt B2} are equivalent.
   If so, returns {\tt TRUE}, otherwise returns {\tt FALSE}.
\endtab
\code

PROCEDURE RandomLattice (B : Basis;   m : BScal;   d : BDim);
\endcode
 \tab Constructs a {\em random\/} basis for an integral lattice,
   chosen with the uniform distribution over the {\em finite\/} set
   of all lattices whose fundamental volume is $m^{d-1}$ and which
   contains $m \Z^d$ as a sublattice, for the specified values of $m$ and $d$.
 \endtab
\code

END LATBASIS.
\endcode

\defmodule {PRIM}

This module provides some procedures to search for integers
$m$ that are prime, for which $r = (m^k-1)/(m-1)$ is also prime
for certain values of $k$, and perhaps for which $(m-1)/2$ is also
prime.  The values of $r$ are always represented as {\tt SuperInteger}'s.

\bigskip\hrule
\code

DEFINITION MODULE PRIM;

FROM SUPINT   IMPORT SuperInteger;
FROM SUPFACT  IMPORT Factors;


PROCEDURE Trier (VAR L : Factors);
\endcode
 \tab 
  Sorts the factors in the list {\tt L} in {\em decreasing\/} order.
 \endtab
\code

PROCEDURE SFindPrimer (k, e : LONGINT;  S1, S2 : SuperInteger;  
                       Safe : BOOLEAN);
\endcode
 \tab
  Finds and prints all integers $m$ such that $S_1 \le m \le S_2$ and    
  such that both $m$ and $r = (m^k-1)/(m-1)$ are prime.   
  If {\tt Safe = TRUE}, it is also required that $(m-1) / 2$ be prime.
  If $k=1$, then $r$ is considered prime.  
  The retained $m$ are printed in decreasing order,
  in the form $m = 2^e \pm x$ and also as decimal integers. 
  The factors of $m-1$ are also printed.
 \endtab
\code

PROCEDURE FindPrimer (k, m0, m1 : LONGINT;  Safe : BOOLEAN);
\endcode
 \tab
  Similar to {\tt SFindPrimemr}, but for values of $m$ less than 
  $2^{31}$.
 \endtab
\code

PROCEDURE FindPrimesmr 
      (k, e, c1, c2 : LONGINT;  Safe : BOOLEAN;   F : ARRAY OF CHAR);
\endcode
 \tab
   Calls {\tt SFindPrimer} for $S_1 = 2^e + c_1$ and $S_2 = 2^e + c_2$, 
   creates the file {\tt F} and prints the results in it.
 \endtab
\code

END PRIM.
\endcode



%
%                   Lower-level modules.
\thispagestyle {nomark}
\addcontentsline{toc}{section}{\protect\numberline{C.}{
  Lower-Level Modules }}
\null\vfill\vfill
\centerline {\large\bf APPENDIX C}
\vfill
\centerline {\large\bf Lower-Level Modules}
\vfill\vfill
%\eject\null\vfill         % Page blanche.

\defmodule {NORM}

This module provides tools to define normalization constants
for the shortest vectors.

\bigskip\hrule
\code

DEFINITION MODULE NORM;

IMPORT BASISLR;
FROM CONFIG IMPORT BDim, MDim;
FROM MULT   IMPORT MScal;

TYPE
   CriterionType = (Spectral, SpectralP, SpectralL1, Beyer, BeyerSpectral,
                    Palpha);
\endcode
 \tabb  Types of merit criteria for lattices.
   {\tt Spectral} means that the criterion is based on the shortest
   vector in the {\em dual\/} lattice.
   In this case, one must specify which approximation is used for $\gamma_t$
   (see {\tt GammaType} below).
   {\tt SpectralP} is similar to {\tt Spectral}, except that
    it is based on the shortest vector in the {\em primal\/} lattice.
   {\tt SpectralL1} is similar to {\tt Spectral}, except that
    the length of the (dual) vectors are measured with the $L_1$ norm,
    minus 1.  The length of the shortest dual vector is then an upper
    bound on the minimal number of hyperplanes that cover all the points
    (\cite{rDIE75a} and \cite{rKNU97a}, Exercises 3.3.4-15 and 16).
   {\tt Beyer} means that the criterion is the Beyer quotient $q_t$.
   {\tt BeyerSpectral} means that both {\tt Beyer} and {\tt Spectral}
   apply simultaneously.
   {\tt Palpha} means that the criterion is $P_\alpha$.
   In this case, one must specify a value of $\alpha$
   and weights may also be used (see the module {\tt PALPHA}).   
 \endtabb
\code

   NormType = (L1Norm, L2Norm, SupNorm);
\endcode
 \tabb  Type of norm used to measure the vector lengths.
  For $V = (v_1,\dots,v_t)$, the $L_1$ norm is 
  $\Vert V\Vert = |v_1|+\cdots+|v_t|$, the $L_2$ (or Euclidean) norm is
  $\Vert V\Vert = (v_1^2+\cdots+v_t^2)^{1/2}$, and the sup norm is
  $\Vert V\Vert = \max(|v_1|,\dots,|v_t|)$.
 \endtabb
\code

   GammaType = (BestLat, Laminated, Minkowski, Rogers, MinkL1);
\endcode
 \tabb  Types of bound on $\gamma_t$ which can be used to normalize 
   the length of the shortest vector.
   The normalized value is the  
   figure of merit used as a selection  criterion.
   {\tt BestLat, Laminated, Minkowski, Rogers} correspond
   respectively to the bounds $\gamma_t^{\rmL}$, $\gamma_t^{\rmB}$, 
   $\gamma_t^{\rmZ}$, and $\gamma_t^{\rmR}$, for the Euclidean norm.
   {\tt MinkL1} must be used for the $\cL_1$ norm and corresponds to 
   $\gamma_t^{\rmM} = (t!)^{1/t}$.
 \endtabb
\code

   Normaliz = POINTER TO InfoNormaliz;
   InfoNormaliz = RECORD
      Name  : ARRAY [0..31] OF CHAR;
      Norm  : NormType;
      GType : GammaType;
      Beta  : LONGREAL;
      Maxt  : BDim;
      Gamma : BASISLR.BVect;
      Cst   : BASISLR.BVect;
      END;
\endcode
  \tabb Data structure used to store a vector of values 
   that are used to normalize the lengths
   $\ell_t$ of the shortest vectors.  
   The array {\tt Gamma} contains values of $\tilde\gamma_t$ that are  
   approximations of $\gamma_t$ which appears in the bounds given
   in Section~1.
   The normalized figure of merit will be $\ell_t/\bar\ell_t$, where
  \eq             
    \bar\ell_t = \cases {
        \tilde\gamma_t m \beta^t       & if $t \le k$;\cr
        \tilde\gamma_t m^{k/t} \beta^t & if $t > k$,\rule{0pt}{14pt}\cr }
                                                    \eqlabel{eq:barellt}
  \endeq
   where $\beta$ is an additional multiplicative factor 
   stored in the variable {\tt Beta} (often equal to 1).
   The vector {\tt Cst} contains the values of $\bar\ell_t^2$ 
   in case of the Euclidean norm, and $\bar\ell_t$ otherwise.

   The structure must be created by the procedure {\tt CreateNormaliz},
   initialized by {\tt InitNormaliz}, and deleted by
   {\tt DeleteNormaliz}.
   The field {\tt Name} is a character string that describes the type of 
   $\tilde\gamma_t$.
   {\tt Norm} is the type of norm that is used to measure
   the vector lengths, {\tt GType} indicates the type of normalization 
   constants $\tilde\gamma_t$, 
   {\tt Beta} is the weighting factor used in (\ref{eq:barellt}),
   {\tt Maxt} is the maximum index for which the arrays have been
   initialized, i.e., the largest value of $t$ that can be considered.
%  Old: With {\tt Critere2 = MinkL1}, {\tt Gamma[t]} contains $t!$.
  \endtabb
\code

PROCEDURE CreateNormaliz (VAR G : Normaliz);
\endcode
 \tab  Creates the structure {\tt G} without initializing anything in it.
 \endtab
\code

PROCEDURE InitNormaliz 
   (   G  : Normaliz;
    Name  : ARRAY OF CHAR;
    Norm  : NormType;
    GType : GammaType;
    Beta  : LONGREAL;
    Maxt  : BDim;
       m  : MScal;
       k  : MDim
   );
\endcode
 \tab  Initializes the structure {\tt G} so that 
   {\tt Gamma[t]} contains the 
   approximation of $\gamma_t$ that corresponds to {\tt T},
   for $1\le t\le {\tt Maxt}$. 
   {\tt Cst[t]} will contain the value of $\bar\ell_t^2$ 
   in case of the Euclidean norm ({\tt Norm = L2Norm}), 
   and $\bar\ell_t$ otherwise,
   where $\bar\ell_t$ is defined in (\ref{eq:barellt}),
   $\tilde\gamma_t$ is determined by {\tt GType},  and
   $\beta$ is the value of {\tt Beta}.

   Normally, this procedure is called with $\beta = 1$.  
   Taking $\beta < 1$ inflates the figure of merit by
   $(1/\beta)^t$, thus weakening the requirements for large $t$ in a 
   worst-case figure of merit such as (\ref{eq:M}).

   The constant {\tt Maxt} cannot exceed 
   the maximum value of $t$ for which the approximation of $\gamma_t$
   has been pre-computed and is available (this depends on the type
   {\tt GType} of $\tilde\gamma_t$).
 \endtab
\code

PROCEDURE DeleteNormaliz (VAR G : Normaliz);
\endcode
 \tab  Deletes the structure {\tt G}.
 \endtab
\code

PROCEDURE WriteNormaliz (G : Normaliz);
\endcode
 \tab  Writes the information contained in {\tt G}.
 \endtab
\code

END NORM.
\endcode


\input CONFIG.tex   \vfill
\input MULT.tex     \vfill
\input BASIS.tex    \vfill
\input CONVERT.tex  \vfill
\defmodule{MULTLI}

This module offers basic tools to manipulate scalars, vectors and matrices
used as multipliers for linear congruential (or multiple recursive) generators.
They are represented using the {\tt LONGINT} type.
Before calling any of the procedures with suffix {\tt Modm}, one must first
call {\tt MSetm} to fix the value of the modulus {\tt m}.
{\tt MaxPeriod} and {\tt PrimElem} can be used to verify the maximal period
conditions for prime moduli given in Knuth \cite{rKNU81a}.
Before calling these procedures, one must first call {\tt SetFactors}
to give or compute the decomposition of $m-1$ and $r$.

\bigskip\hrule
\code

DEFINITION MODULE MULTLI;

FROM SUPINT  IMPORT SuperInteger;
FROM SUPFACT IMPORT Factors;
FROM CONFIG  IMPORT MDim;

TYPE
   MScal   = LONGINT;
   MVect   = ARRAY MDim OF MScal;
   MMatr   = ARRAY MDim OF MVect;

VAR
   MZero, MOne, MTwo : MScal;            (* Constants 0, 1, and 2.       *)

PROCEDURE MCreateScal (VAR S : MScal);
\endcode
  \tab  Creates {\tt S}. (Is there for compatibility with {\tt MULTSI}.
  Does nothing in the case of {\tt MULTLI}.) \endtab
\code

PROCEDURE MDeleteScal (VAR S : MScal);
\endcode
  \tab Deallocates the memory allocated to {\tt S} by a previous call to {\tt
  MCreateScal}.\endtab
\code

PROCEDURE MCreateVect (VAR V : MVect; k1, k2 : MDim);
\endcode
  \tab  Creates entries of {\tt V}, using {\tt MCreateScal},
  for indices {\tt k1}  to  {\tt k2} inclusive.
  If  {\tt k2} is smaller than  {\tt k1}  then does nothing.\endtab
\code

PROCEDURE MDeleteVect (VAR V : MVect; k1, k2 : MDim);
\endcode
  \tab  Deallocates the memory  that was allocated for the creation of V,
  of type {\tt MVect}, at a previous call to {\tt MCreateVect} with indices 
  from {\tt k1}  to  {\tt k2} inclusive.
  \endtab
\code

PROCEDURE MCreateMatr (VAR M : MMatr; k1, k2 : MDim);
\endcode
  \tab  Creates the elements of {\tt M}, as in {\tt MCreateVect}.\endtab
\code

PROCEDURE MDeleteMatr (VAR M : MMatr; k1, k2 : MDim);
\endcode
  \tab Deallocates the memory  that was allocated for the creation of M,
  of type {\tt MMatr}, at a previous call to {\tt MCreateMatr} with parameters 
  {\tt k1} and {\tt k2}.
  \endtab
\code

PROCEDURE MReadScal (VAR S : MScal);
\endcode
  \tab  Reads {\tt S} from current intput file.
  \endtab
\code

PROCEDURE MWriteScal (S : MScal);
\endcode
  \tab  Writes {\tt S} to current output file.
  \endtab
\code

PROCEDURE MWriteScalBin (S : MScal);
\endcode
  \tab  Writes {\tt S} to current output file in binary form.
  \endtab
\code

PROCEDURE MWriteVect (VAR V : MVect;  n : MDim);
\endcode
  \tab  Writes content of {\tt V} in current output file.
  \endtab
\code

PROCEDURE MScalToLI (S : MScal) : LONGINT;
\endcode
  \tab  Returns the {\tt LONGINT} representation of {\tt S}.\endtab
\code

PROCEDURE MScalToLR (S : MScal) : LONGREAL;
\endcode
  \tab  Returns the {\tt LONGREAL} representation of {\tt S}.\endtab
\code

PROCEDURE MScalToSup (S : MScal;  T : SuperInteger);
\endcode
  \tab  Converts {\tt S} to the {\tt SuperInteger T} (which must have been
   created before).\endtab
\code

PROCEDURE SupToMScal (T : SuperInteger;  VAR S : MScal);
\endcode
  \tab  Converts {\tt T} to the {\tt MScal S} (which must have been
   created before).\endtab
\code

PROCEDURE LIToMScal (I : LONGINT;  VAR S : MScal);
\endcode
  \tab  Converts {\tt I} to the {\tt MScal S}.\endtab
\code

PROCEDURE MAffect (S : MScal;  VAR T : MScal);
\endcode
  \tab  Gives to {\tt T} the value of {\tt S}. \endtab
\code

PROCEDURE MEqual (S, T : MScal): BOOLEAN;
\endcode
  \tab  Returns {\tt TRUE} if {\tt S} is equal to {\tt T},
  {\tt FALSE} otherwise.
  \endtab
\code

PROCEDURE MGreater (S1, S2 : MScal): BOOLEAN;
\endcode
  \tab  Returns {\tt TRUE} if {\tt S1} is greater than {\tt S2},
  {\tt FALSE} otherwise.
  \endtab
\code

PROCEDURE MSmaller (S1, S2 : MScal): BOOLEAN;
\endcode
  \tab  Returns {\tt TRUE} if {\tt S1} is smaller than {\tt S2},
  {\tt FALSE} otherwise.
  \endtab
\code

PROCEDURE MChangeSign (VAR S : MScal);
\endcode
  \tab  Changes the sign of S. \endtab
\code

PROCEDURE MInc (VAR S : MScal; I : LONGINT);
\endcode
  \tab  Adds I to S.  The result is in S.  \endtab
\code

PROCEDURE MAdd (S1, S2 : MScal;   VAR T : MScal);
\endcode
  \tab  Adds {\tt S1} to {\tt S2} and puts the results in {\tt T}.\endtab
\code

PROCEDURE MSubtract (S1, S2 : MScal;   VAR T : MScal);
\endcode
  \tab  Subtracts {\tt S2} from {\tt S1} and puts the results in {\tt T}.\endtab
\code

PROCEDURE MAffectIdentMatr (VAR M : MMatr;  n : MDim);
\endcode
  \tab  Initializes {\tt M[1..n,1..n]} to the identity matrix.\endtab
\code

PROCEDURE MCopyVect (VAR U, V : MVect;  n : MDim);
\endcode
  \tab  Copies vector {\tt U}, of dimension {\tt n}, into vector {\tt V}.
  \endtab
\code

PROCEDURE MCopyMatr (VAR M, N : MMatr;   n : MDim);
\endcode
  \tab  Copies matrix {\tt M}, of dimension {\tt n$\times$n}, into matrix {\tt N}.
  \endtab
\code

PROCEDURE MModifVect (VAR U,V : MVect;   R : MScal; n : MDim);
\endcode
  \tab  Adds {\tt R} times the vector {\tt V} to the vector {\tt U}.
  \endtab
\code

PROCEDURE MPowerLI (S : MScal;  i : LONGINT;  VAR R : MScal);
\endcode
   \tab  Returns $S^i$ in {\tt R}. \endtab
\code

PROCEDURE MMultiply (T, S : MScal;   VAR R : MScal);
\endcode
   \tab  Returns the product of {\tt T} by {\tt S} in {\tt R}.
   \endtab
\code

PROCEDURE MQuotient (T, S : MScal;   VAR Q : MScal);
\endcode
   \tab  Returns the (truncated) quotient of {\tt T} by {\tt S} in {\tt Q}.
   \endtab
\code

PROCEDURE MModulo (T, S : MScal;   VAR R : MScal);
\endcode
   \tab  Returns {\tt T} modulo {\tt S} in {\tt R}.
   \endtab
\code

PROCEDURE MSetm (m : MScal);
\endcode
  \tab  Sets the value of the modulus {\tt m} which will be used from now on
  in the procedures below.\endtab
\code

PROCEDURE MGetm (VAR m : MScal);
\endcode
  \tab  Returns in {\tt m} the value of the current modulus.\endtab
\code

PROCEDURE MProdModm (S, T : MScal;  VAR R : MScal);
\endcode
   \tab Returns in {\tt R} the product of {\tt S} by {\tt T} modulo {\tt m}.
    Assumes that {\tt S} and {\tt T} are between {\tt -m} and {\tt m}.
    \endtab
\code

PROCEDURE MProdScalModm (VAR U, V : MVect;   VAR Q, R : MScal;   n : MDim);
\endcode
   \tab Returns in {\tt R} the scalar product of {\tt U} by {\tt V}
    modulo {\tt m}, using components {\tt 1}
    to {\tt n} of these vectors.  {\tt Q} contains the quotient.\endtab
\code

PROCEDURE MModifVectModm (VAR U, V : MVect;  R : MScal;  n : MDim);
\endcode
  \tab  Adds {\tt R} times the vector {\tt V} to the vector {\tt U}.
  \endtab
\code

PROCEDURE MPowerPolyMod (VAR A : MVect;  n : MDim;  q : MScal;  VAR R : MVect);
\endcode
   \tab Returns in {\tt R} the coefficients of the polynomial
     $x^q \mod f(x) ~(\mod m)$,
     where $f(x) = x^n - A_1 x^{n-1} -\cdots- A_n$.
   \endtab
\code

PROCEDURE SetFactors (VAR ListmMS1, Listr : Factors;   k : MDim);
\endcode
   \tab Sets the factors of $m-1$ and $r ~[= (m^k-1)/(m-1)]$. If
        {\tt ListmMS1} is {\tt NIL}, then the procedure will calculate and
        return the factors of $m-1$. If not, then the procedure will take the
        list of factors (see {\tt SUPFACT} of {\tt SENTIERS}) and verify if
        they are factors of $m-1$. The program will stop if the product does
        not give $m-1$. The same things are done with {\tt Listr} and
        the factors of $r$.
   \endtab
\code

PROCEDURE PrimElem (VAR A : MVect) : BOOLEAN;
\endcode
   \tab  Verifies a first condition for $f(x)$ to be primitive modulo $m$,
    namely that $(-1)^{k-1} a_k$ is a primitive element modulo $m$.
    Must have called {\tt SetFactors} before.
   \endtab
\code

PROCEDURE MaxPeriod (VAR A : MVect) : BOOLEAN;
\endcode
   \tab  Verifies the remaining conditions for $f(x)$ to be primitive modulo
    $m$, assuming that {\tt PrimElem} has returned {\tt TRUE}.
   \endtab
\code

END MULTLI.
\endcode

\defmodule{MULTSI}

Same as {\tt MULTLI}, except that the values are
represented using the {\tt SuperInteger} type, and that {\tt MCreateScal},
{\tt MCreateVect}, and {\tt MCreateMatr} are required to create the 
{\tt SuperInteger}s.  The initial memory space (in words) allocated 
for the digits of the {\tt SuperInteger}s thus created is 2 to the power
{\tt MSizeSI} (currently 4 words). (See {\tt Create} in {\tt SUPINT}).  
\bigskip
\hrule
\code

DEFINITION MODULE MULTSI;

FROM SUPINT  IMPORT SuperInteger;
FROM SUPFACT IMPORT Factors;
FROM CONFIG  IMPORT MMaxDim, MDim;

TYPE
   MScal   = SuperInteger;
   MVect   = ARRAY MDim OF MScal;
   MMatr   = ARRAY MDim OF MVect;


VAR
   MZero, MOne, MTwo : MScal;          (* These should be constants      *)


PROCEDURE MCreateScal (VAR S : MScal);
PROCEDURE MDeleteScal (VAR S : MScal);
PROCEDURE MCreateVect (VAR V : MVect; k1, k2 : MDim);
PROCEDURE MDeleteVect (VAR V : MVect; k1, k2 : MDim);
PROCEDURE MCreateMatr (VAR M : MMatr; k1, k2 : MDim);
PROCEDURE MDeleteMatr (VAR M : MMatr; k1, k2 : MDim);
PROCEDURE MReadScal (VAR S : MScal);
PROCEDURE MEqual (S, T : MScal): BOOLEAN;
PROCEDURE MGreater (S1, S2 : MScal): BOOLEAN;
PROCEDURE MSmaller (S1, S2 : MScal): BOOLEAN;
PROCEDURE MChangeSign (S : MScal);
PROCEDURE MInc (VAR S : MScal; I : LONGINT);
PROCEDURE MAdd (S1, S2 : MScal; VAR S3 : MScal);
PROCEDURE MSubtract (S1, S2 : MScal;  VAR S3 : MScal);
PROCEDURE MWriteScal (S : MScal);
PROCEDURE MWriteScalBin (S : MScal);
PROCEDURE MWriteVect (VAR V : MVect;  n : MDim);
PROCEDURE MScalToLI (S : MScal) : LONGINT;
PROCEDURE MScalToLR (S : MScal) : LONGREAL;
PROCEDURE MScalToSup (S: MScal; T : SuperInteger);
PROCEDURE LIToMScal (I : LONGINT;  VAR S : MScal);
PROCEDURE MAffect (S : MScal;  VAR T : MScal);
PROCEDURE MAffectIdentMatr (VAR M : MMatr;  n : MDim);
PROCEDURE MCopyVect (VAR U, V : MVect;  n : MDim);
PROCEDURE MCopyMatr (VAR M, N : MMatr; n : MDim);
PROCEDURE MModifVect (VAR U,V : MVect; R : MScal; n : MDim);
PROCEDURE MPowerLI (S : MScal;  i : LONGINT;  VAR R : MScal);
PROCEDURE MMultiply (T, S : MScal;   VAR  R : MScal);
PROCEDURE MQuotient (T, S : MScal;   VAR  Q : MScal);
PROCEDURE MModulo (T, S : MScal;   VAR  R : MScal);
PROCEDURE MSetm (m : MScal);
\endcode 
\code
PROCEDURE MGetm (VAR m : MScal);
PROCEDURE MProdModm (S, T : MScal;  VAR R : MScal);
PROCEDURE MProdScalModm (VAR U, V : MVect;  VAR a, R : MScal;   n : MDim);
\endcode 
\code
PROCEDURE MModifVectModm (VAR U,V : MVect; R : MScal; n : MDim);
PROCEDURE MPowerPolyMod (VAR A : MVect; n : MDim; q : MScal; VAR R : MVect);
PROCEDURE SetFactors (VAR ListmMS1, Listr : Factors; k : MDim);
\endcode
\code
PROCEDURE PrimElem (VAR A : MVect) : BOOLEAN;
PROCEDURE MaxPeriod (VAR A : MVect) : BOOLEAN;
END MULTSI.
\endcode

\defmodule{BASISLR}

This module offers basic tools to manipulate scalars, vectors, and matrices
that are used to represent lattice bases, using the {\tt LONGREAL} type.
It is used by the modules {\tt LATBASIS}, {\tt REDBAS}, etc.,
when this representation type is used.
See that module for more details.

\bigskip\hrule
\code

DEFINITION MODULE BASISLR;

FROM SUPINT   IMPORT SuperInteger;
FROM CONFIG   IMPORT BMaxDim, BDim;

TYPE
   BScal   = LONGREAL;
   BVect   = ARRAY BDim OF BScal;
   BMatr   = ARRAY BDim OF BVect;
   VectLI  = ARRAY BDim OF LONGINT;

VAR
   BZero, BOne, BMinusOne : BScal;    (* 0, 1, and -1.                   *)

PROCEDURE BCreateScal (VAR S : BScal);

PROCEDURE BDeleteScal (VAR S : BScal);

PROCEDURE BCreateVect (VAR V : BVect; d1, d2 : BDim);

PROCEDURE BDeleteVect (VAR V : BVect; d1, d2 : BDim);

PROCEDURE BCreateMatr (VAR M : BMatr; d1, d2 : BDim);

PROCEDURE BDeleteMatr (VAR M : BMatr; d1, d2 : BDim);
\endcode
  \tab These 6 procedures do nothing.
  There are there only for compatibility with the corresponding
  procedures in {\tt BASISSI}.
  \endtab
\code

PROCEDURE BReadScal (VAR S : BScal);

PROCEDURE BWriteScal (S : BScal;   nb : CARDINAL);
\endcode
  \tab  Writes content of {\tt S} in current output file.
        Uses a total of at least {\tt nb} positions.
  \endtab
\code

PROCEDURE BWriteVect (VAR V : BVect;  n : BDim;  nb : CARDINAL);
\endcode
  \tab  Writes content of {\tt V} for indices {\tt [1..n]}
        in current output file.
        Uses a total of at least {\tt nb} positions for each number.
  \endtab
\code

PROCEDURE BWriteMatr (VAR M : BMatr;  n : BDim; nb : CARDINAL);
\endcode
  \tab  Writes content of {\tt M} for indices {\tt [1..n][1..n]}
        in current output file.
        Uses a total of at least {\tt nb} positions for each number.
  \endtab
\code

PROCEDURE BProdScalVect (s : BScal;  VAR V : BVect;  d : BDim);
\endcode
  \tab Multiplies vector {\tt V} by scalar {\tt s}. \endtab
\code

PROCEDURE BProdScal (VAR U, V : BVect;  n : BDim;  VAR S : BScal);
\endcode
   \tab Returns in {\tt S} the scalar product of {\tt U} by {\tt V},
    using components from 1 to {\tt n}. \endtab
\code

PROCEDURE BProdScalLR (VAR U, V : BVect;  n : BDim) : LONGREAL;
\endcode
   \tab Returns the scalar product of {\tt U} by {\tt V},
    using components from 1 to {\tt n}. \endtab
\code

PROCEDURE BProdScalSI (VAR U, V : BVect; n : BDim; Res : SuperInteger);
\endcode
  \tab  Returns in {\tt S} the scalar product of {\tt U} by {\tt V},
    using components from 1 to {\tt n}.  {\tt U} and {\tt V} are
    rounded and converted into {\tt SuperInteger} representation before
    the computation of the scalar product. \endtab
\code

PROCEDURE BL1NormLR (VAR V : BVect;  n : BDim) : LONGREAL;
\endcode
   \tab Returns the $L_1$ norm of {\tt V},
    using components from 1 to {\tt n}. \endtab
\code

PROCEDURE BL2NormLR (VAR V : BVect;  n : BDim) : LONGREAL;
\endcode
   \tab Returns the $L_2$ norm of {\tt V},
    using components from 1 to {\tt n}. \endtab
\code

PROCEDURE BL1NormSI (VAR V : BVect;  n : BDim;  S1 : SuperInteger);
\endcode
   \tab Returns in {\tt S1} the $L_1$ norm of {\tt V},
    using components from 1 to {\tt n}. \endtab
\code

PROCEDURE BSupNormLR (VAR V : BVect;  n : BDim) : LONGREAL;
\endcode
   \tab Returns the sup norm of {\tt V},
    using components from 1 to {\tt n}. \endtab
\code

PROCEDURE BSupNormSI (VAR V : BVect;  n : BDim;  S1 : SuperInteger);
\endcode
   \tab Returns in {\tt S1} the sup norm of {\tt V},
    using components from 1 to {\tt n}. \endtab
\code

PROCEDURE BModifVectLR (VAR U, V : BVect;  c : LONGREAL;  n : BDim);
\endcode
  \tab  Adds {\tt c} times the vector {\tt V} to the vector {\tt U}.
   Assumes that both vectors have dimension {\tt n}.\endtab
\code

PROCEDURE BModifVectLI (VAR U, V : BVect;  c : LONGINT;  n : BDim);
\endcode
  \tab  Adds {\tt c} times the vector {\tt V} to the vector {\tt U}.
   Assumes that both vectors have dimension {\tt n}.\endtab
\code

PROCEDURE BModifVect (VAR U, V : BVect;  c : BScal;  n : BDim);
\endcode
  \tab  Adds {\tt c} times the vector {\tt V} to the vector {\tt U}.
   Assumes that both vectors have dimension {\tt n}.\endtab
\code

PROCEDURE BChangeSign (VAR S : BScal);
\endcode
  \tab  Changes the sign of S. \endtab
\code

PROCEDURE BChangeSignVect (VAR V : BVect;  n : BDim);
\endcode
  \tab  Changes the sign of elements {\tt 1} to {\tt n} of vector {\tt V}.
  \endtab
\code

PROCEDURE BAffect (S : BScal;  VAR T : BScal);
\endcode
  \tab  Gives to {\tt T} the value of {\tt S}.\endtab
\code

PROCEDURE BExchange (VAR S, T : BScal);

PROCEDURE BEqual (S, T : BScal): BOOLEAN;
\endcode
  \tab  Returns {\tt TRUE} if {\tt S} is equal to {\tt T},
   {\tt FALSE} otherwise.
  \endtab
\code

PROCEDURE BSmaller (S, T : BScal): BOOLEAN;
\endcode
  \tab  Returns {\tt TRUE} if {\tt S} is smaller than {\tt T},
  {\tt FALSE} otherwise.
  \endtab
\code

PROCEDURE BGreater (S, T : BScal): BOOLEAN;
\endcode
  \tab  Returns {\tt TRUE} if {\tt S} is greater than {\tt T},
  {\tt FALSE} otherwise.
  \endtab
\code

PROCEDURE BAbsSizeOrd (S, T : BScal): INTEGER;
\endcode
  \tab  Returns -1 (resp. 0, +1) according as {\tt S} $>$ (resp. $=, <$) {\tt T}. 
  \endtab
\code

PROCEDURE BAbs (S : BScal;  VAR T : BScal);
\endcode
  \tab  Returns in {\tt T} the absolute value of {\tt S}. \endtab
\code

PROCEDURE BSetNeg (VAR S : BScal);
\endcode
  \tab  Changes the sign of {\tt S} if it is not negative.
  \endtab
\code

PROCEDURE BInc (VAR S : BScal; i : LONGINT);
\endcode
  \tab  Adds  {\tt i}  to  {\tt S}.
  \endtab
\code

PROCEDURE BAdd (S, T : BScal;  VAR Res : BScal);
\endcode
  \tab  Returns in {\tt Res} the sum of {\tt S} and {\tt T}.  \endtab
\code

PROCEDURE BMultiply (S, T : BScal;   VAR Res : BScal);
\endcode
  \tab  Returns in {\tt Res} the product of {\tt S} by {\tt T}. \endtab
\code

PROCEDURE BDivideRound (S, T : BScal; VAR Res : BScal);
\endcode
  \tab  Returns in {\tt Res} the result of the division of {\tt S}
        by {\tt T}.  The result is rounded.  For results just in the
        middle of two integers, the value returned by the function is
        the integer the nearest of zero.
        Ex.: If $S=5.0$ and $T=10.0$, then the returned value is 0.
            If $S=-5.0$ and $T=10.0$, then the returned value is also 0.
  \endtab
\code

PROCEDURE BCopyVect (VAR U, V : BVect;  n : BDim);
\endcode
  \tab  Copies the vector {\tt U}, of dimension {\tt n}, into vector {\tt V}.
  \endtab
\code

PROCEDURE BCopyMatr (VAR M1, M2 : BMatr;  n : BDim);
\endcode
  \tab  Copies matrix {\tt M1}, of dimension {\tt n}, into matrix {\tt M2}.
  \endtab
\code

PROCEDURE BConvLR (S : BScal): LONGREAL;
\endcode
  \tab  Returns the {\tt LONGREAL} representation of {\tt S}.
  Is there only for compatibility with {\tt BASISSI}. \endtab
\code

PROCEDURE BConvSI (R : BScal;  S : SuperInteger);
\endcode
  \tab  Returns in {\tt S} the {\tt SuperInteger} representation of {\tt R}.
  \endtab
\code

PROCEDURE LIToBScal (i : LONGINT; VAR S : BScal);
\endcode
   \tab Converts i to  {\tt S}\endtab
\code

PROCEDURE BAffectIdentMatr (VAR M : BMatr;  n : BDim);
\endcode
  \tab  Initializes {\tt M[1..n,1..n]} to the identity matrix.\endtab
\code

PROCEDURE BTrMatr (M : BMatr;   VAR Mtr : BMatr;   d : BDim);
\endcode
  \tab  Computes the transposed matrix of {\tt M}, and
        returns the result in {\tt Mtr}.
  \endtab
\code

PROCEDURE BCanonicalVect (VAR V : BVect;  i : BDim);
\endcode
   \tab Produces the vector with all coordinates equal to zero with the
   exception of the $i$-th coordinate which is equal to one.
   \endtab
\code

PROCEDURE BRandomScal (VAR s : BScal; m : BScal);
\endcode
   \tab Produces a random integer in the interval $[0,m)$, with a uniform
   distribution. \endtab
\code

PROCEDURE GCD2vect (VAR V : VectLI;   k, n : BDim) : LONGINT;
\endcode
 \tab  Computes and returns the greatest common divisor of the vectors
  {\tt V[k],\dots,V[n]}.
 \endtab
\code

PROCEDURE BEuclide (A, B : BScal;  VAR C, D, E, F, G : BScal);
\endcode
  \tab For given $A$ and $B$, returns $C$, $D$, $E$, $F$, $G$ such that: 
   \begin {eqnarray*}
     C  A + D  B  &=& G = \mbox{gcd} (A,B) \\
     E  A + F  B  &=& 0,
   \end {eqnarray*}
where {gcd}  is the greatest common denominator.
  \endtab
\code

END BASISLR.
\endcode

\defmodule{BASISSI}

Same as {\tt BASISLR}, except that the values are
represented using the {\tt SuperInteger} type.

\bigskip\hrule
\code

DEFINITION MODULE BASISSI;

FROM SUPINT   IMPORT SuperInteger;
FROM CONFIG   IMPORT BDim;

CONST
   BSizeSI = 2;                        (* Used in BCreateScal.            *)
TYPE
   BScal   = SuperInteger;
   BVect   = ARRAY BDim OF BScal;
   BMatr   = ARRAY BDim OF BVect;

VAR
   BZero, BOne, BMinusOne : BScal;     (* 0, 1, and -1                    *)

PROCEDURE BCreateScal (VAR S : BScal);
\endcode
  \tab  Creates the {\tt SuperInteger S}.
   The initial memory space (in words)
   allocated for the digits of {\tt S} is 2 to the
   power {\tt BSizeSI} (currently 4 words).
   (See {\tt Create} in {\tt SUPINT}).\endtab
\code

PROCEDURE BDeleteScal (VAR S : BScal);
\endcode
  \tab  Deallocates the memory allocated for the {\tt SuperInteger S},
  at a previous call to {\tt BCreateScal}. \endtab
\code

PROCEDURE BCreateVect (VAR V : BVect; d1, d2 : BDim);
\endcode
  \tab  Creates the {\tt SuperInteger} entries for indices 
  {\tt d1}  to  {\tt d2} inclusive
  (i.e., extends {\tt V[0..d1-1]} to {\tt V[0..d2]}).
  If  {\tt d2} is smaller than  {\tt d1}  then does nothing. \endtab
\code

PROCEDURE BDeleteVect (VAR V : BVect; d1, d2 : BDim);
\endcode
  \tab  Deallocates the memory allocated to the {\tt SuperInteger}
  elements of {\tt V[d1..d2]}.
  \endtab
\code

PROCEDURE BCreateMatr (VAR M : BMatr; d1, d2 : BDim);
\endcode
  \tab  Creates the {\tt SuperInteger} entries for indices 
   {\tt d1}  to  {\tt d2} inclusive; i.e., extends the square matrix 
   {\tt M[0..d1-1, 0..d1-1]} to {\tt M[0..d2, 0..d2]}.
   If {\tt d1} $> 0$, creates the missing entries in the {\tt d1} 
   vectors already there, in addition to creating the {\tt d2 - d1 + 1}
   missing vectors.
\endtab
\code

PROCEDURE BDeleteMatr (VAR M : BMatr; d1, d2 : BDim);
\endcode
  \tab  Does exactly the opposite of {\tt BCreateMatr (M, d1, d2)}.
  \endtab
\code

PROCEDURE BReadScal (VAR S : BScal);

PROCEDURE BWriteScal (S : BScal;   nb : CARDINAL);
\endcode
  \tab  Writes the contents of {\tt S} to the current output.
        Uses {\tt nb} digits per line.  \endtab
\code

PROCEDURE BWriteVect (VAR V : BVect;  n : BDim;  nb : CARDINAL);
\endcode
  \tab  Writes contents of {\tt V} to current output.
        Uses {\tt nb} digits per line for each number.  \endtab
\code

PROCEDURE BWriteMatr (VAR M : BMatr;  n : BDim;  nb : CARDINAL);
\endcode
  \tab  Writes contents of {\tt M} to current output.
        Uses nb digits per line for each number.  \endtab
\code

PROCEDURE BProdScalVect (s : BScal;  VAR V : BVect;  d : BDim);
PROCEDURE BProdScal (VAR U, V : BVect;  n : BDim;  VAR S : BScal);
PROCEDURE BProdScalLR (VAR U, V : BVect;  n : BDim) : LONGREAL;
PROCEDURE BProdScalSI (VAR U, V : BVect; n : BDim; Res : SuperInteger);
\endcode
\code
PROCEDURE BL1NormLR (VAR V : BVect;  n : BDim) : LONGREAL;
PROCEDURE BL1NormSI (VAR V : BVect;  n : BDim;  S1 : SuperInteger);
PROCEDURE BL2NormLR (VAR V : BVect;  n : BDim) : LONGREAL;
PROCEDURE BSupNormLR (VAR V : BVect;  n : BDim) : LONGREAL;
PROCEDURE BSupNormSI (VAR V : BVect;  n : BDim;  S1 : SuperInteger);
\endcode
\code
PROCEDURE BModifVectLR (VAR U, V : BVect;  c : LONGREAL;  n : BDim);
PROCEDURE BModifVectLI (VAR U, V : BVect;  c : LONGINT;  n : BDim);
PROCEDURE BModifVect (VAR U, V : BVect;  c : BScal;  n : BDim);
\endcode
\code
PROCEDURE BChangeSign (VAR S : BScal);
PROCEDURE BChangeSignVect (VAR V : BVect;  n : BDim);
PROCEDURE BAffect (S : BScal;  VAR T : BScal);
PROCEDURE BExchange (VAR S, T : BScal);
\endcode
\code
PROCEDURE BEqual (S, T : BScal): BOOLEAN;
PROCEDURE BSmaller (S, T : BScal): BOOLEAN;
PROCEDURE BGreater (S, T : BScal): BOOLEAN;
\endcode
\code
PROCEDURE BAbsSizeOrd (S, T : BScal): INTEGER;
PROCEDURE BAbs (S : BScal; VAR T : BScal);
PROCEDURE BSetNeg (VAR S : BScal);
PROCEDURE BInc (VAR S : BScal; i : LONGINT);
PROCEDURE BAdd (S, T : BScal; VAR Res : BScal);
PROCEDURE BMultiply (S, T : BScal; VAR Res : BScal);
PROCEDURE BDivideRound (S, T : BScal; VAR Res : BScal);
PROCEDURE BCopyVect (VAR U, V : BVect;  n : BDim);
PROCEDURE BCopyMatr (VAR M1, M2 : BMatr;  n : BDim);
PROCEDURE BConvLR (S : BScal): LONGREAL;
PROCEDURE BConvSI (R : BScal; S : SuperInteger);
PROCEDURE LIToBScal (i : LONGINT; VAR S : BScal);
PROCEDURE BAffectIdentMatr (VAR M : BMatr;  n : BDim);
PROCEDURE BTrMatr (M : BMatr; VAR Mtr : BMatr; d : BDim);
PROCEDURE BCanonicalVect (VAR V : BVect; i : BDim);
PROCEDURE BRandomScal (VAR s : BScal; m : BScal);

PROCEDURE BEuclide (A, B : BScal;  VAR C, D, E, F, G : BScal);
\endcode
  \tab For given $A$ and $B$, returns $C$, $D$, $E$, $F$, $G$ such that: 
   \begin {eqnarray*}
     C  A + D  B  &=& G = \mbox{gcd} (A,B) \\
     E  A + F  B  &=& 0.
   \end {eqnarray*}
where {gcd}  is the greatest common denominator.
  \endtab
\code

END BASISSI.
\endcode

\defmodule{LILR}

Procedures to convert between {\tt LONGINT} (multipliers) and
{\tt LONGREAL} (elements of basis vectors).
This is one of the four versions of the module {\tt CONVERT}.
The other versions are {\tt LISI}, {\tt SILR}, {\tt SISI}.
Which one is used depends on the current representation, as selected by
the command {\tt select}.

\bigskip\hrule
\code

DEFINITION MODULE LILR;

FROM MULTLI  IMPORT MScal;             (* LONGINT.                        *)
FROM BASISLR IMPORT BScal;             (* LONGREAL.                       *)
FROM SUPINT  IMPORT SuperInteger;

PROCEDURE MScalToBScal (S : MScal;  VAR R : BScal);
\endcode
  \tab  Converts {\tt S} to {\tt R}
   (from {\tt LONGINT} to {\tt LONGREAL}).\endtab
\code

PROCEDURE BScalToMScal (S : BScal;  VAR R : MScal);
\endcode
  \tab  Converts {\tt S} to {\tt R}
   (from {\tt LONGREAL} to {\tt LONGINT}).\endtab
\code

PROCEDURE SupToBScal (S : SuperInteger;  VAR R : BScal);
\endcode
  \tab Converts {\tt S} to {\tt R}
   (from SuperInteger to {\tt LONGREAL}).\endtab
\code

PROCEDURE SupToMScal (S : SuperInteger;  VAR R : MScal);
\endcode
  \tab Converts {\tt S} to {\tt R}
   (from SuperInteger to {\tt LONGINT}).\endtab
\code

PROCEDURE BScalToSup (S : BScal;  R : SuperInteger);
\endcode
  \tab Converts {\tt S} to {\tt R}
   (from {\tt LONGREAL} to SuperInteger).\endtab
\code

PROCEDURE MScalToSup (S : MScal;  R : SuperInteger);
\endcode
  \tab Converts {\tt S} to {\tt R}
   (from {\tt LONGINT} to SuperInteger).\endtab
\code

END LILR.
\endcode

\defmodule {LISI}

Procedures to convert between {\tt LONGINT} (multipliers) and
{\tt SuperInteger} (elements of basis vectors).
Similar to {\tt LILR}.

\bigskip\hrule
\code

DEFINITION MODULE LISI;

FROM MULTLI  IMPORT MScal;             (* LONGINT.                        *)
FROM BASISSI IMPORT BScal;             (* SuperInteger.                   *)
FROM SUPINT  IMPORT SuperInteger;

PROCEDURE MScalToBScal (S : MScal;  VAR R : BScal);
\endcode
  \tab  Converts {\tt S} to {\tt R}
   (from {\tt LONGINT} to {\tt SuperInteger}).\endtab
\code

PROCEDURE BScalToMScal (S : BScal;  VAR R : MScal);
\endcode
  \tab  Converts {\tt S} to {\tt R}
   (from {\tt SuperInteger} to {\tt LONGINT}).\endtab
\code

PROCEDURE SupToBScal (S : SuperInteger;  VAR R : BScal);
\endcode
  \tab Copies {\tt S} to {\tt R}.\endtab
\code

PROCEDURE SupToMScal (S : SuperInteger;  VAR R : MScal);
\endcode
  \tab Converts {\tt S} to {\tt R}
   (from SuperInteger to {\tt LONGINT}).\endtab
\code

PROCEDURE BScalToSup (S : BScal;  R : SuperInteger);
\endcode
  \tab Copies {\tt S} to {\tt R}.\endtab
\code

PROCEDURE MScalToSup (S : MScal;  R : SuperInteger);
\endcode
  \tab Converts {\tt S} to {\tt R}
   (from {\tt LONGINT} to SuperInteger).\endtab
\code

END LISI.
\endcode

\defmodule{SILR}

Procedures to convert between {\tt SuperInteger} (multipliers) and
{\tt LONGREAL} (elements of basis vectors).
Similar to {\tt LILR}.

\bigskip\hrule
\code

DEFINITION MODULE SILR;

FROM SUPINT  IMPORT SuperInteger;
FROM MULTSI  IMPORT MScal;             (* SuperInteger.                   *)
FROM BASISLR IMPORT BScal;             (* LONGREAL.                       *)

PROCEDURE MScalToBScal (S : MScal;   VAR R : BScal);
\endcode
  \tab  Converts {\tt S} to {\tt R}
   (from {\tt SuperInteger} to {\tt LONGREAL}).\endtab
\code

PROCEDURE BScalToMScal (S : BScal;   VAR R : MScal);
\endcode
  \tab  Converts {\tt S} to {\tt R}
   (from {\tt LONGREAL} to {\tt SuperInteger}).\endtab
\code

PROCEDURE SupToBScal (S : SuperInteger;   VAR R : BScal);
\endcode
  \tab Converts {\tt S} to {\tt R}
   (from SuperInteger to {\tt LONGREAL}).\endtab
\code

PROCEDURE SupToMScal (S : SuperInteger;   VAR R : MScal);
\endcode
  \tab Copies {\tt S} to {\tt R}.\endtab
\code

PROCEDURE BScalToSup (S : BScal;   R : SuperInteger);
\endcode
  \tab Converts {\tt S} to {\tt R}
   (from {\tt LONGREAL} to SuperInteger).\endtab
\code

PROCEDURE MScalToSup (S : MScal;   R : SuperInteger);
\endcode
  \tab Copies {\tt S} to {\tt R}.\endtab
\code

END SILR.
\endcode

\defmodule{SISI}

Procedures to convert between {\tt SuperInteger} (multipliers) and
{\tt SuperInteger} (elements of basis vectors).
Simply copies the value (equivalent to {\tt Affect} in {\tt SUPINT}).

\bigskip\hrule
\code

DEFINITION MODULE SISI;

FROM MULTSI  IMPORT MScal;             (* SuperInteger.                   *)
FROM BASISSI IMPORT BScal;             (* SuperInteger.                   *)
FROM SUPINT  IMPORT SuperInteger;

PROCEDURE MScalToBScal (S : MScal;  VAR R : BScal);
\endcode
  \tab  Copies {\tt S} to {\tt R}.\endtab
\code

PROCEDURE BScalToMScal (S : BScal;  VAR R : MScal);
\endcode
  \tab  Copies {\tt S} to {\tt R}.\endtab
\code

PROCEDURE SupToBScal (S : SuperInteger; VAR R : BScal);
\endcode
  \tab Copies {\tt S} to {\tt R}.\endtab
\code

PROCEDURE SupToMScal (S : SuperInteger; VAR R : MScal);
\endcode
  \tab Copies {\tt S} to {\tt R}.\endtab
\code

PROCEDURE BScalToSup (S : BScal; R : SuperInteger);
\endcode
  \tab Copies {\tt S} to {\tt R}.\endtab
\code

PROCEDURE MScalToSup (S : MScal; R : SuperInteger);
\endcode
  \tab Copies {\tt S} to {\tt R}.\endtab
\code

END SISI.
\endcode

\defmodule {select}

This command will change automatically the representations of the
moduli and multipliers in the modules {\tt MULT}, {\tt BASIS} and
{\tt CONVERT}. To run this program, type ``{\tt select}'' $xx$ $yy$, where
$xx$ is {\tt LI} or {\tt SI}, and $yy$ is {\tt LR} or {\tt SI}.
Parameters $xx$ and $yy$  determine the types used for multipliers and for 
the moduli, respectively.
The program copies {\tt MULT}$xx$ in {\tt MULT}, {\tt BASIS}$yy$ in
{\tt BASIS}, and $xxyy$ in {\tt CONVERT}. 
The module names are also changed appropriately.
%  Under VAX/VMS, this program will recompile all the modules of
%  {\tt ExamLCG} that are affected by the change. 
%  Under MS-DOS, this will be done automatically by the $TopSpeed$ 
%  make facility.

%
%%%%%%%%%%%%%%%%%%%%%%%%%%%%%%%%%%%%%%%%%%%%%%%%%%%%%%%%%%%%%%%%%%%%%%%%%%%%
\thanks

This work has been supported by NSERC-Canada grant \# ODGP0110050
and FCAR-Qu\'ebec grant \# 93ER1654 to the first author.
Fran\c cois Blouin, Anna Bragina, Ajmal Chaumun, Marco Jacques, 
Fran\c cois Paradis, and Jos\'ee Turgeon 
have participated in the implementation of this software package,
since 1988.
We thank L. Afflerbach and H. Grothe for helpful discussions 
in the early stage of the development, in 1989.

%%%%%%%%%%%%%%%%%%%%%%%%%%%%%%%%%%%%%%%%%%%%%%%%%%%%%%%%%%%%%%%%%%%%%%%%%%%%

\bibliography {random,simul,math,ift,vrt}
\bibliographystyle {plain}

\end {document}
